 % =============================================================================
% Signals PROJECT BLACKBOOK (B.E. THESIS REPORT)
% A. P. Shah Institute of Technology, Thane -- University of Mumbai
% Academic Year: 2025-2026
%
% Compile with: pdflatex main.tex  (run twice for ToC)
% =============================================================================

\documentclass[12pt,a4paper,oneside,openany]{book}

% ---- Fonts (Times New Roman equivalent, pre-installed in all TeX distros) -
\usepackage{mathptmx}
\usepackage{tocloft}

% ---- Page geometry per APSIT guidelines ----------------------------------
\usepackage[
  a4paper,
  top=30mm,
  bottom=25mm,
  left=30mm,
  right=20mm,
  headheight=16pt,
  headsep=10mm,
  footskip=12mm
]{geometry}

% ---- Line spacing (1.5) --------------------------------------------------
\usepackage{setspace}
\onehalfspacing

% ---- Basic packages ------------------------------------------------------
\usepackage{graphicx}
\graphicspath{{figures/}{./}}
\usepackage{xcolor}
\usepackage{array}
\usepackage{booktabs}
\usepackage{tabularx}
\usepackage{longtable}
\usepackage{multirow}
\usepackage{makecell}
\usepackage{amsmath,amssymb}
\usepackage{enumitem}
\usepackage{titlesec}
\usepackage{fancyhdr}
\usepackage{caption}
\usepackage{float}
\usepackage{url}
\usepackage[hidelinks]{hyperref}
\usepackage{etoolbox}
\usepackage{ragged2e}

% ---- Paragraph indent ----------------------------------------------------
\setlength{\parindent}{12mm}
\setlength{\parskip}{4pt}


% ---- TOC / LOF / LOT title formatting -------------------------------------
\renewcommand{\cfttoctitlefont}{\hfill\fontsize{18}{22}\selectfont\bfseries}
\renewcommand{\cftaftertoctitle}{\hfill}

\renewcommand{\cftloftitlefont}{\hfill\fontsize{18}{22}\selectfont\bfseries}
\renewcommand{\cftafterloftitle}{\hfill}

\renewcommand{\cftlottitlefont}{\hfill\fontsize{18}{22}\selectfont\bfseries}
\renewcommand{\cftafterlottitle}{\hfill}

% Move titles upward
\setlength{\cftbeforetoctitleskip}{-20pt}
\setlength{\cftaftertoctitleskip}{20pt}

\setlength{\cftbeforeloftitleskip}{-20pt}
\setlength{\cftafterloftitleskip}{20pt}

\setlength{\cftbeforelottitleskip}{-20pt}
\setlength{\cftafterlottitleskip}{20pt}

% ---- Captions ------------------------------------------------------------
\captionsetup{labelfont=bf,labelsep=endash,font=small}

% ---- Chapter formatting --------------------------------------------------
% Chapter: "CHAPTER N" centered 18pt, then chapter title 18pt bold centered
\titleformat{\chapter}[display]
  {\normalfont\centering\bfseries}
  {\fontsize{18}{22}\selectfont CHAPTER \thechapter}
  {12mm}
  {\fontsize{18}{22}\selectfont}
\titlespacing*{\chapter}{0pt}{60mm}{15mm}

% Centered format for unnumbered chapters (Contents, List of Figures, etc.)
\titleformat{name=\chapter,numberless}[block]
  {\normalfont\centering\bfseries}
  {}
  {0pt}
  {\fontsize{18}{22}\selectfont}
\titlespacing*{name=\chapter,numberless}{0pt}{20mm}{15mm}

% Section: 16pt bold left-aligned with 15mm space above and below
\titleformat{\section}
  {\normalfont\fontsize{16}{20}\bfseries}
  {\thesection}{0.6em}{}
\titlespacing*{\section}{0pt}{15mm}{8mm}

% Sub-section: 14pt bold left-aligned
\titleformat{\subsection}
  {\normalfont\fontsize{14}{18}\bfseries}
  {\thesubsection}{0.6em}{}
\titlespacing*{\subsection}{0pt}{10mm}{6mm}

% Sub-sub-section: 12pt bold
\titleformat{\subsubsection}
  {\normalfont\fontsize{12}{15}\bfseries}
  {\thesubsubsection}{0.5em}{}
\titlespacing*{\subsubsection}{0pt}{6mm}{4mm}

% ---- Header / footer style ----------------------------------------------
\pagestyle{fancy}
\fancyhf{}
\renewcommand{\headrulewidth}{0pt}
\renewcommand{\footrulewidth}{0pt}
\fancyhead[L]{\fontsize{10}{12}\selectfont\itshape Chapter \thechapter\quad\leftmark}
\fancyfoot[R]{\thepage}
\renewcommand{\chaptermark}[1]{\markboth{#1}{}}

\fancypagestyle{plain}{%
  \fancyhf{}%
  \fancyfoot[R]{\thepage}%
  \renewcommand{\headrulewidth}{0pt}%
}

% Remove page number from first page of each chapter (per guidelines)
\patchcmd{\chapter}{\thispagestyle{plain}}{\thispagestyle{empty}}{}{}

% ---- Placeholder figure command -----------------------------------------
% Creates a bordered box acting as a figure placeholder if image is missing.
\newcommand{\figurebox}[2][10cm]{%
  \fbox{\parbox[c][#1][c]{0.92\linewidth}{\centering\itshape #2}}%
}

\newcommand{\boxfig}[3]{%
  \begin{figure}[H]
    \centering
    \figurebox[#1]{#2}
    \caption{#3}
  \end{figure}
}

% ---- Table of contents depth --------------------------------------------
\setcounter{tocdepth}{2}
\setcounter{secnumdepth}{3}

% =============================================================================
\begin{document}
% =============================================================================

% ########## FRONT MATTER ##########
\frontmatter
\pagenumbering{roman}
\pagestyle{empty}

% -------------- TITLE PAGE (single page, singlespaced) ---------------------
\begin{titlepage}
\thispagestyle{empty}
\begin{singlespace}
\begin{center}
    % \vspace*{2mm}
    {\Large \textbf{Signals: FORECASTING MEMECOIN PRICES\\
    BASED ON MARKET Signals, RUG CHECKS, AND\\
    TWITTER SENTIMENT DATA} \par}

    \vspace{1cm}

    {\fontsize{12}{14}\selectfont Submitted in partial fulfilment of the requirements of the degree of}\\[6mm]

    {\fontsize{14}{17}\selectfont\bfseries BACHELOR OF COMPUTER ENGINEERING}\\[8mm]

    {\fontsize{12}{14}\selectfont by}\\[4mm]

    {\fontsize{13}{16}\selectfont Om Bhojane, 22102097}\\[2mm]
    \vspace*{2mm}
    {\fontsize{13}{16}\selectfont Mihir Amin, 22102088}\\[2mm]
    \vspace*{2mm}
    {\fontsize{13}{16}\selectfont Sairaj Bokand, 22102133}\\[2mm]
    \vspace*{2mm}
    {\fontsize{13}{16}\selectfont Shreyas Bhalekar, 22102168}\\[8mm]

    {\fontsize{13}{16}\selectfont Guide}\\[2mm]
    \vspace*{2mm}
    {\fontsize{13}{16}\selectfont Dr. Uttam D. Kolekar\\[6mm]

    \includegraphics[width=48mm]{apsit.png}\\[6mm]

    {\fontsize{13}{16}\selectfont Department of Computer Engineering}\\[2mm]
    \vspace*{2mm}
    {\fontsize{13}{16}\selectfont\bfseries A. P. SHAH INSTITUTE OF TECHNOLOGY, THANE}\\[2mm]
    \vspace*{2mm}
    {\fontsize{12}{15}\selectfont University of Mumbai}\\[2mm]
    \vspace*{2mm}
    {\fontsize{12}{15}\selectfont 2025--2026}
\end{center}
\end{singlespace}
\end{titlepage}

% -------------- CERTIFICATE -------------------------------------------------
\clearpage
\thispagestyle{empty}
\begin{singlespace}
\noindent
\begin{minipage}[c]{15mm}
  \includegraphics[width=30mm]{apsit.png}
\end{minipage}%
\hfill
\begin{minipage}[c]{\dimexpr\linewidth-32mm\relax}
  \centering
  {\fontsize{15}{18}\selectfont\bfseries A. P. SHAH INSTITUTE OF TECHNOLOGY, THANE}
\end{minipage}
% \vspace{10mm}

\begin{center}
  {\fontsize{20}{24}\selectfont\bfseries CERTIFICATE}\\[15mm]
\end{center}
\begin{spacing}{1.5}
{\setstretch{1.3}
\noindent This is to certify that the project entitled
\textbf{``Signals: Forecasting Memecoin Prices Based on Market Signals, Rug Checks, and Twitter Sentiment Data''} is a bonafide work of \textbf{Om Bhojane (22102097)}, \textbf{Mihir Amin (22102088)}, \textbf{Sairaj Bokand (22102133)}, and \textbf{Shreyas Bhalekar (22102168)} submitted to the University of Mumbai in partial fulfilment of the requirement for the award of the degree of \textbf{Bachelor of Engineering in Computer Engineering}.
\par}
\end{spacing}
\end{singlespace}
\vspace{4cm}

\noindent
\begin{minipage}{0.45\textwidth}
    \raggedright
    \rule{4cm}{0.4pt} \\   % signature line (left aligned now)
    Guide \\
    Dr. Uttam D. Kolekar
\end{minipage}
\hfill
\begin{minipage}{0.45\textwidth}
    \raggedright
    \rule{4cm}{0.4pt} \\
    Project Coordinator \\
    Dr. Rushikesh R. Nikam \\
\end{minipage}

\vspace{4cm}

\noindent
\begin{minipage}{0.45\textwidth}
    \raggedright
    \rule{4cm}{0.4pt} \\
    Head of Department \\
    Dr. Sachin H. Malave \\
\end{minipage}
\hfill
\begin{minipage}{0.45\textwidth}
    \raggedright
    \rule{4cm}{0.4pt} \\
    Principal \\
    Dr. Uttam D. Kolekar \\
\end{minipage}

% \noindent Date:\hspace{10mm}\\[2mm]
% Place: Thane

% -------------- PROJECT REPORT APPROVAL (CERTIFICATE STYLE MATCH) -----------
\clearpage
\thispagestyle{empty}

\begin{singlespace}

\noindent
\begin{minipage}[c]{15mm}
  \includegraphics[width=30mm]{apsit.png}
\end{minipage}%
\hfill
\begin{minipage}[c]{\dimexpr\linewidth-32mm\relax}
  \centering
  {\fontsize{15}{18}\selectfont\bfseries A. P. SHAH INSTITUTE OF TECHNOLOGY, THANE}
\end{minipage}

\vspace{8mm}

\begin{center}
  {\fontsize{20}{24}\selectfont\bfseries Project Report Approval for B.E.}\\[12mm]
\end{center}

\begin{spacing}{1.5}
{\setstretch{1.3}
\noindent
This project report entitled
\textbf{``Signals: Forecasting Memecoin Prices Based on Market Signals, Rug Checks, and Twitter Sentiment Data''}
by \textbf{Om Bhojane, Mihir Amin, Sairaj Bokand, and Shreyas Bhalekar}
is approved for the degree of \textbf{Bachelor of Engineering in Computer Engineering}.
\textbf{2025--2026}.
\par}
\end{spacing}

\end{singlespace}

\vspace{20mm}

% ---------------- Signature Section ----------------

\noindent
\begin{tabular}{p{8cm} p{5cm}}
\textbf{Examiner Name} & \textbf{Signature} \\
\end{tabular}

\vspace{10mm}

\noindent
\begin{tabular}{p{8cm} p{5cm}}
1.\ \rule{6.5cm}{0.4pt} & \rule{4.5cm}{0.4pt} \\[12mm]
2.\ \rule{6.5cm}{0.4pt} & \rule{4.5cm}{0.4pt} \\
\end{tabular}

\vspace{18mm}

\noindent
\textbf{Date:} \\[6mm]

\noindent
\textbf{Place: Thane}

% -------------- DECLARATION (FIXED EXACT MATCH) ----------------
\clearpage
\thispagestyle{empty}

\begin{center}
  {\fontsize{16}{18}\selectfont\bfseries Declaration}\\[10mm]
\end{center}

\noindent
{\fontsize{11}{14}\selectfont
We declare that this written submission represents our ideas in our own words and where others' ideas or words have been included, we have adequately cited and referenced the original sources. We also declare that we have adhered to all principles of academic honesty and integrity and have not misrepresented or fabricated or falsified any idea/data/fact/source in our submission. We understand that any violation of the above will be cause for disciplinary action by the Institute and can also evoke penal action from the sources which have thus not been properly cited or from whom proper permission has not been taken when needed.
}

\vspace{20mm}

% ---------------- RIGHT-ALIGNED SIGNATURE BLOCK ----------------

\begin{flushright}

\begin{tabular}{l}

\rule{5.5cm}{0.4pt} \\
\vspace{1cm}
Om Bhojane - 22102097 \\[12mm]

\rule{5.5cm}{0.4pt} \\
\vspace{0.8cm}
Mihir Amin - 22102088 \\[12mm]
% \vspace{2cm}

\rule{5.5cm}{0.4pt} \\
\vspace{0.8cm}
Sairaj Bokand - 22102133 \\[12mm]
% \vspace{2cm}

\rule{5.5cm}{0.4pt} \\
Shreyas Bhalekar - 22102168 \\

\end{tabular}

\end{flushright}

\vspace{15mm}

\noindent

% -------------- ABSTRACT ----------------------------------------------------
\clearpage
\thispagestyle{empty}
\begin{center}
  {\fontsize{18}{22}\selectfont\bfseries Abstract}\\[10mm]
\end{center}

\noindent
Memecoin markets exhibit extreme volatility, sentiment-driven price swings, and a high prevalence of rug-pull exploits, making them fundamentally incompatible with traditional cryptocurrency analytics. Existing tools operate in silos: price-prediction models rely solely on technical indicators, rug-pull detectors analyse only smart-contract features, and sentiment analysers aggregate social-media mentions without integrating them into actionable trading decisions. In this project, an integrated multi-agent platform called \textbf{Signals} was developed to forecast memecoin prices and assess token risk in real time by unifying market Signals, on-chain rug checks, and Twitter sentiment data. The system orchestrates specialised LLM-based agents---a Market Agent, a Rug-Check Agent, a Social Agent, and a Predictor Agent---through a ReAct-style reasoning loop with episodic and semantic memory, enabling continual self-improvement across market regimes. On-chain safety analysis covers liquidity locks, holder concentration, mint authority, and contract-level red flags, while the sentiment pipeline ingests Twitter activity and reconciles it against on-chain truth. A custom synthetic market simulator with slippage and regime modelling provides a zero-risk evaluation environment for the agentic trader. The platform delivers transparent, interpretable verdicts with sub-second latency, demonstrating the viability of unified, self-improving, multi-signal analysis for the extreme volatility of modern memecoin markets.

\vspace{8mm}

\noindent
\textit{\textbf{Keywords:} Memecoin Forecasting, Rug-Pull Detection, Multi-Agent LLM, Twitter Sentiment Analysis, On-Chain Analytics, ReAct Agents, Episodic Memory, Solana, Decentralised Finance, Agentic Trading}

% -------------- CONTENTS ----------------------------------------------------
\clearpage
\pagestyle{empty}
\thispagestyle{empty}
\begin{center}
  {\fontsize{18}{22}\selectfont\bfseries Contents}
\end{center}
\vspace{8mm}
\makeatletter
\@starttoc{toc}
\makeatother
\thispagestyle{empty}

\clearpage
\thispagestyle{empty}
\begin{center}
  {\fontsize{18}{22}\selectfont\bfseries List of Figures}
\end{center}
\vspace{8mm}
\makeatletter
\@starttoc{lof}
\makeatother
\thispagestyle{empty}

\clearpage
\thispagestyle{empty}
\begin{center}
  {\fontsize{16}{18}\selectfont\bfseries Abbreviations}\\[8mm]
\end{center}

\begin{center}
\renewcommand{\arraystretch}{1.2}
\setlength{\tabcolsep}{7pt}
\small

\begin{tabular}{|>{\centering\arraybackslash}p{2.2cm}|>{\centering\arraybackslash}p{8.5cm}|}
\hline
\textbf{LLM}  & Large Language Model \\ \hline
\textbf{RL}   & Reinforcement Learning \\ \hline
\textbf{DeFi} & Decentralised Finance \\ \hline
\textbf{DEX}  & Decentralised Exchange \\ \hline
\textbf{AMM}  & Automated Market Maker \\ \hline
\textbf{API}  & Application Programming Interface \\ \hline
\textbf{ReAct} & Reasoning + Acting (agent framework) \\ \hline
\textbf{CoT}  & Chain-of-Thought (prompting) \\ \hline
\textbf{RSI}  & Relative Strength Index \\ \hline
\textbf{MACD} & Moving Average Convergence Divergence \\ \hline
\textbf{BB}   & Bollinger Bands \\ \hline
\textbf{GBM}  & Geometric Brownian Motion \\ \hline
\textbf{BSC}  & Binance Smart Chain \\ \hline
\textbf{KYC}  & Know Your Customer \\ \hline
\textbf{NLP}  & Natural Language Processing \\ \hline
\textbf{LSTM} & Long Short-Term Memory \\ \hline
\textbf{TFT}  & Temporal Fusion Transformer \\ \hline
\textbf{JSON} & JavaScript Object Notation \\ \hline
\textbf{UML}  & Unified Modelling Language \\ \hline
\textbf{DFD}  & Data Flow Diagram \\ \hline
\textbf{SDG}  & Sustainable Development Goal \\ \hline
\textbf{P\&L} & Profit and Loss \\ \hline
\textbf{SLM}  & Small Language Model \\ \hline
\textbf{CNN}  & Convolutional Neural Network \\ \hline
\textbf{DQN}  & Deep Q-Network \\ \hline
\textbf{SAC}  & Soft Actor-Critic \\ \hline
\textbf{DDPG} & Deep Deterministic Policy Gradient \\ \hline
\textbf{MoE}  & Mixture of Experts \\ \hline
\textbf{CEX}  & Centralised Exchange \\ \hline
\textbf{BCI}  & Brain-Computer Interface \\ \hline
\end{tabular}
\end{center}

% ########## MAIN MATTER ##########
\mainmatter
\pagestyle{fancy}
\pagenumbering{arabic}
\setcounter{page}{1}

% =============================================================================
% CHAPTER 1 -- INTRODUCTION
% =============================================================================
\chapter{Introduction}
\label{ch:intro}

The global cryptocurrency market has experienced an unprecedented structural shift over the past three years, driven not by technological utility but by the rise of a new and highly speculative asset class known as the \emph{memecoin}. Unlike traditional cryptocurrencies whose valuations are anchored in consensus mechanisms, use-cases, or network throughput, memecoins derive their perceived value from internet culture, community virality, and the gravitational pull of social-media sentiment. In 2024 and 2025, platforms such as Pump.fun single-handedly accelerated this phenomenon by allowing anyone to launch a new Solana token in less than sixty seconds with zero capital, resulting in the generation of an estimated 71\% of all new Solana tokens during this period~\cite{ref10}. Yet, despite this staggering throughput, fewer than two per cent of launched tokens ``graduate'' to a sustainable level of on-chain liquidity on decentralised exchanges such as Raydium or Orca. The remaining 98\% are abandoned, manipulated, or fraudulently drained by their creators within hours or days.

The price volatility of memecoins is unlike anything traditional finance has ever encountered. The MURICA token, analysed as one of the live case studies in this project, recorded a 3{,}576\% price change within a single 24-hour window. Such extreme movements are frequently accompanied by sophisticated exploits collectively known as \emph{rug pulls}---attacks in which developers deliberately design, launch, and hype a token only to withdraw liquidity, dump their allocation on retail buyers, and disappear. A recent survey of the literature identified 35 distinct rug-pull attack patterns~\cite{ref20}, and Chainalysis' 2025 Crypto Crime Report estimates that rug pulls and exit scams collectively stole more than \$2.6 billion from retail investors over the past twelve months. These exploits are compounded by the fact that existing analytical infrastructure remains fundamentally fragmented: price-prediction models rely solely on technical indicators~\cite{ref4,ref7}, rug-pull detectors analyse only smart-contract features~\cite{ref5}, and sentiment analysis tools merely aggregate social-media mentions and hashtags without integrating them into any actionable trading decision~\cite{ref6}.

\section{Motivation}
The central problem facing a modern memecoin trader is therefore one of \emph{signal integration}. A tool that excels at detecting rug pulls is useless if it cannot tell the user whether the token is trending on Twitter; a sentiment analyser that correctly identifies viral hype is dangerous if it ignores the fact that the liquidity pool is unlocked and the top ten wallets control 72\% of the supply. Recent multi-agent LLM trading systems such as TradingAgents~\cite{ref23}, FinCon~\cite{ref24}, CryptoTrade~\cite{ref25} and FinMem~\cite{ref26} have demonstrated the potential of LLM-based financial reasoning, yet none of these systems addresses the unique challenges of memecoin markets: extreme volatility, prevalent frauds, and social-sentiment-driven price dynamics operating on minute-level timescales. This gap motivated the design of a unified system capable of fusing market, safety, and social Signals into a single coherent trading decision.

\section{Problem Definition}
Formally, given a token $\tau$ on a blockchain $\mathcal{B}$ with a feature vector $x \in \mathbb{R}^{23}$ comprising technical, on-chain and social Signals, the objective is to produce a trading decision $d \in \{\mathrm{BUY}, \mathrm{SELL}, \mathrm{HOLD}\}$ together with a confidence value $c \in [0,1]$ and a categorical risk assessment $r \in \{\mathrm{LOW}, \mathrm{MEDIUM}, \mathrm{HIGH}\}$, subject to an end-to-end latency constraint $t_{\mathrm{response}} < 30$~seconds. This formulation distinguishes the present work from traditional cryptocurrency-prediction research by demanding \emph{simultaneous} fraud detection, sentiment integration and real-time decision making under severe temporal constraints.

\section{Project Contributions}
The Signals project delivers five primary contributions:

\begin{enumerate}[leftmargin=12mm]
\item \textbf{Tri-pillar signal fusion:} Engineering of twenty-three distinct features spanning eleven technical indicators, seven on-chain safety metrics, and five social sentiment Signals, analysed by four specialised LLM agents with a composite scoring framework (40\% technical + 40\% on-chain + 20\% social) and cross-domain safety overrides.

\item \textbf{Multi-agent ReAct architecture:} Implementation of a five-phase orchestrated THINK-ACT-OBSERVE-ACT-REFLECT pipeline with Chain-of-Thought reasoning, tool use, and confidence-score aggregation across four role-specialised agents running at distinct temperature profiles.

\item \textbf{Agentic self-improving trader:} Introduction of a THINK-ACT-REFLECT-LEARN trading loop backed by a dual memory system (episodic trading journal + semantic rule base) that enables autonomous natural-language rule discovery without gradient-based learning, inspired by Tulving's episodic-semantic memory distinction~\cite{ref45}.

\item \textbf{Open-source LLM benchmark:} A head-to-head benchmark of eight open-source and proprietary LLMs on crypto trading tasks, finding that open-source models (Grok-4.20 at +34.59\%, DeepSeek-V3.1 at +7.30\%) consistently outperform proprietary alternatives (Gemini-3-Pro at $-10.94\%$, Claude-Sonnet-4.5 at $-19.15\%$).

\item \textbf{Synthetic market simulator:} A regime-switching Geometric Brownian Motion simulator with jump diffusion, correlated on-chain and social metric generation, and realistic technical-indicator computation, enabling controlled and repeatable agent evaluation.
\end{enumerate}

\section{Significance of the Work}
By integrating previously siloed analytical domains into a unified agentic intelligence platform, Signals establishes a comprehensive defence, monitoring, and decision-support layer for retail investors operating within the highly volatile and speculative memecoin ecosystem. Unlike traditional tools that focus on isolated Signals such as price trends or social sentiment, Signals fuses multi-modal data streams—including on-chain activity, social media dynamics, developer behavior, and liquidity patterns—into a coherent analytical pipeline. This holistic approach enables early detection of anomalous patterns such as coordinated hype cycles, liquidity manipulation, and potential rug-pull events, thereby significantly improving risk awareness for non-expert participants.

In addition, Signals contributes to the broader discourse on agentic AI systems by demonstrating how modular, collaborative agents can be orchestrated to solve complex, real-world problems. The architecture highlights the practical viability of multi-agent coordination, tool usage, and memory augmentation in building scalable and adaptive intelligent systems. As such, this work not only advances the state-of-the-art in crypto analytics but also serves as a reference design for next-generation AI-driven financial infrastructures.

\section{Organisation of the Report}
The remainder of this report is structured to provide a systematic progression from foundational concepts to implementation and evaluation. Chapter~\ref{ch:lit} presents a comprehensive survey of related work spanning six key research domains: cryptocurrency price prediction, memecoin-specific analysis, rug-pull detection mechanisms, multi-agent LLM applications in finance, agentic AI frameworks, and sentiment analysis techniques tailored to crypto markets. This chapter establishes the theoretical and technological context for the proposed system.

Chapter~\ref{ch:limit} critically examines the limitations of existing approaches, highlighting challenges such as lack of interpretability, fragmented analysis pipelines, poor generalization in volatile markets, and absence of standardized evaluation frameworks. Based on this analysis, the chapter identifies the specific research gaps that motivate the development of Signals.

Chapter~\ref{ch:prob} formalizes the problem definition, clearly outlining the objectives, assumptions, and scope of the project. It also defines key performance metrics and success criteria used throughout the study. Chapter~\ref{ch:proposed} provides an in-depth description of the proposed system architecture, including detailed UML diagrams, agent interactions, data flow pipelines, and core algorithms. The methodology section elaborates on the design choices, model configurations, and reasoning mechanisms employed.

Chapter~\ref{ch:setup} describes the experimental setup, including the software stack, hardware infrastructure, datasets, and evaluation protocols used to validate the system. Chapter~\ref{ch:impl} presents the implementation details, followed by experimental results, comparative analysis, case studies, and ablation experiments to assess the contribution of individual components.

Chapter~\ref{ch:plan} outlines the project timeline through a detailed Gantt chart, illustrating the phased development and milestones achieved. Finally, Chapter~\ref{ch:conclusion} summarizes the key findings and contributions of the work, while Chapter~\ref{ch:future} discusses potential extensions, scalability considerations, and directions for future research in agentic AI and financial intelligence systems.

% =============================================================================
% CHAPTER 2 -- LITERATURE SURVEY
% =============================================================================
\chapter{Literature Survey}
\label{ch:lit}

The design of Signals draws on research from six distinct but interlocking areas: cryptocurrency price prediction, memecoin market analysis, rug-pull detection, multi-agent LLM systems for finance, agentic AI frameworks, and sentiment analysis for crypto markets. This chapter presents a critical appraisal of the state of the art in each of these six areas and identifies the gaps that motivated the present work.

\section{Cryptocurrency Price Prediction}
Researchers have investigated cryptocurrency price prediction for several years using a wide variety of statistical, machine-learning, and deep-learning techniques. Early work in this domain relied on traditional econometric models such as ARIMA and GARCH, which attempted to model temporal dependencies and volatility clustering in financial time series. However, these models often struggled with the highly non-linear, noisy, and non-stationary nature of cryptocurrency markets.

With the evolution of machine learning, more robust predictive models emerged. Wang \textit{et al.}~\cite{ref7} provide a comprehensive review of deep learning approaches for cryptocurrency price prediction and highlight that architectures such as Long Short-Term Memory (LSTM) networks and transformer-based models are widely adopted for forecasting tasks. LSTMs are particularly effective due to their ability to capture long-term dependencies in sequential data, while transformers leverage attention mechanisms to identify relevant temporal patterns across varying time horizons.

Traditional machine-learning methods such as XGBoost~\cite{ref4} have also demonstrated strong predictive performance. These models are especially effective in short-horizon forecasting scenarios, where market behavior tends to be relatively stable and decision boundaries are easier to learn. Their ability to incorporate engineered features such as technical indicators, trading volume, and price momentum makes them highly competitive against deep learning approaches in certain settings.

More advanced sequence models such as the Temporal Fusion Transformer (TFT)~\cite{ref21} have further improved forecasting capabilities. TFT integrates static features (e.g., asset-specific characteristics) with time-varying inputs (e.g., historical prices, exogenous Signals) and produces calibrated probabilistic forecasts rather than single-point estimates. This probabilistic output is particularly valuable in financial contexts, where uncertainty quantification plays a critical role in risk management and decision-making. Additionally, TFT supports multi-horizon forecasting, enabling predictions across multiple future time steps simultaneously.

In parallel, recent research has begun to explore the application of large language models (LLMs) in finance. Benchmarks such as FinBen~\cite{ref37} evaluate LLM performance across a variety of financial tasks, including price prediction, sentiment analysis, and numerical reasoning. These models are capable of processing unstructured data sources such as news articles, financial reports, and social media content factors that significantly influence cryptocurrency markets. Domain-specific foundation models such as BloombergGPT~\cite{ref38} and FinGPT~\cite{ref39} further demonstrate that finance-specialised pre-training can substantially enhance downstream task performance, particularly in tasks that require contextual understanding of financial language and events.

Despite these advancements, nearly all of the existing research focuses on well-established cryptocurrencies such as Bitcoin and Ethereum, which exhibit relatively mature market dynamics, higher liquidity, and stronger institutional participation. These characteristics make them more amenable to traditional modeling approaches.

In contrast, memecoins represent a fundamentally different and less-explored category of digital assets. Their price movements are often driven by sudden hype cycles, viral social media trends, influencer endorsements, and community-led coordination rather than intrinsic value or fundamental indicators. Additionally, memecoins are highly susceptible to abrupt events such as regulatory announcements and malicious activities like rug pulls, where project developers abandon the asset after artificially inflating its price.

These unique characteristics introduce extreme volatility, regime shifts, and unpredictable behavioral patterns that are difficult for conventional models to capture. Furthermore, the lack of structured financial data and reliable historical trends adds another layer of complexity. As a result, the dynamics of memecoin markets remain substantially under-explored in current prediction literature, highlighting a significant research gap and the need for hybrid modeling approaches that combine time-series analysis, sentiment extraction, and adaptive learning frameworks.

\section{Memecoin Market Analysis}
A small but rapidly growing body of work has begun to investigate the unique behaviour of memecoin markets, which differ significantly from traditional cryptocurrency ecosystems in both structure and driving forces. MemeChain~\cite{ref1} introduces a multi-agent simulation framework to model how memecoin markets evolve over time, capturing interactions between heterogeneous participants such as early adopters, speculators, and opportunistic traders. This approach highlights how decentralized and loosely coordinated communities can collectively influence price dynamics, often resulting in emergent behaviours such as sudden rallies or collapses.

ME2F~\cite{ref2} further explores the fragility of memecoin ecosystems by analysing their market microstructures. The study identifies recurring patterns that precede instability events, such as liquidity imbalances, abrupt trading volume spikes, and asymmetric participation. These findings suggest that memecoin markets are inherently unstable due to their reliance on speculative sentiment rather than intrinsic value, making them highly sensitive to external triggers.

Recent studies focusing on the Pump.fun platform~\cite{ref10} provide deeper insight into early-stage token dynamics. These works demonstrate that multiple modalities of information available at the time of token launch including token names, descriptive text, initial liquidity parameters, and even early visual chart patterns carry predictive Signals regarding a token’s long-term survival. Notably, linguistic cues (e.g., humor, cultural references) and visual trends (e.g., early price momentum patterns) can act as early indicators of community engagement and subsequent adoption.

Building on these ideas, MemeTrans~\cite{ref11} explores transformer-based architectures tailored specifically for memecoin price prediction. By leveraging attention mechanisms, these models can capture complex temporal dependencies and rapidly changing market Signals, which are characteristic of memecoin environments. Similarly, CoinCLIP~\cite{ref33} extends prediction capabilities into the multimodal domain by linking visual meme content with market behaviour. Using contrastive vision-language learning, CoinCLIP learns joint embeddings of meme images and their corresponding price trajectories, demonstrating that visual culture memes, virality, and symbolism plays a measurable role in influencing financial outcomes.

Parallel research on manipulation detection~\cite{ref34} investigates how coordinated pump-and-dump schemes artificially inflate memecoin prices. These studies reveal that such schemes often involve synchronized activity across multiple social platforms, including rapid dissemination of promotional content, influencer endorsements, and coordinated buying behaviour. Additionally, research on meme-coin copy-trading behaviour~\cite{ref35} shows that investors frequently mimic the actions of perceived “successful” traders, creating tight feedback loops. This leads to strong herd behaviour, where price movements are amplified not by fundamentals but by collective imitation and fear of missing out (FOMO).

Despite these valuable contributions, most existing studies remain limited in scope. Many focus on offline analysis using historical datasets or examine only a single dimension of the ecosystem such as price prediction, sentiment analysis, or manipulation detection in isolation. Crucially, there is a lack of integrated systems that combine real-time data ingestion, multimodal signal processing (text, visuals, market data), safety-metric evaluation (e.g., rug-pull risk, liquidity health), and price prediction within a unified pipeline.

This gap presents a significant opportunity for future research. Developing a holistic, real-time analytical framework that captures the full complexity of memecoin markets could enable more robust predictions, early risk detection, and better decision-making for participants operating in these highly volatile environments.

\section{Rug-Pull Detection}
Detecting rug-pull scams through smart-contract analysis has become an important area of research. Several systems have shown strong results in identifying malicious token contracts. CRPWarner~\cite{ref5} reports a detection accuracy of 91.8\%, while RPHunter~\cite{ref12} achieves 95.3\% by combining behavioural and static features. TokenScout~\cite{ref13} reaches approximately 98\% accuracy by evaluating live token behaviour and contract characteristics across multiple dimensions.

Other approaches study on-chain transaction patterns directly. TM-RugPull~\cite{ref36} analyses transaction motifs to identify common activity patterns that appear immediately before a rug-pull event, while a broader systematisation-of-knowledge study~\cite{ref20} identifies 35 distinct categories of rug-pull attacks and demonstrates that many existing detection tools recognise only 25--63\% of known patterns. SolRPDS~\cite{ref3} specifically focuses on the Solana blockchain and provides a dedicated dataset and detection pipeline for Solana rug pulls.

Although these systems are effective at flagging suspicious smart contracts, they typically operate as standalone tools and do not combine contract-risk analysis with other factors such as price prediction or social sentiment. Consequently, they do not provide the holistic risk perspective that memecoin traders need when evaluating new tokens in the few minutes after launch.

\section{Multi-Agent LLM Systems for Finance}
The use of multi-agent LLM systems in financial analysis has recently attracted substantial attention. TradingAgents~\cite{ref23} uses a debate-style framework in which multiple agents discuss and evaluate stock-trading decisions before reaching a collective conclusion. FinCon~\cite{ref24} introduces a synthesised multi-agent system that uses conceptual verbal reinforcement to enhance portfolio management decisions. In the cryptocurrency space, CryptoTrade~\cite{ref25} uses a reflective LLM-based agent to guide zero-shot cryptocurrency trading, but crucially it does not include any safety checks related to on-chain fraud risk.

FinMem~\cite{ref26} proposes a memory-augmented LLM system for financial decision making that stores and retrieves past experiences through a layered memory structure, while FinAgent~\cite{ref27} demonstrates how multimodal LLM agents can be used for trading by combining heterogeneous data sources and integrating external tools. Work on multi-source information fusion~\cite{ref17} shows that integrating multiple data streams supports more accurate crypto market analysis than any single source alone. M-DQN~\cite{ref14} integrates sentiment into deep Q-learning for Bitcoin trading, and complementary work applies SAC and DDPG to cryptocurrency portfolio optimisation~\cite{ref15}, while broader surveys of reinforcement learning in quantitative finance~\cite{ref16} catalogue the methodological landscape.

Although these studies demonstrate the potential of LLM-based agents in finance, most of them focus on traditional financial assets or generic cryptocurrency trading. Essentially no prior work has applied multi-agent LLM reasoning specifically to memecoin markets, and existing systems typically do not include rug-pull detection as part of the trading decision process---an omission that makes them structurally unsuitable for the memecoin setting.

\section{Agentic AI Frameworks}
Recent research on agentic AI has introduced several frameworks that allow language models to perform more complex reasoning and acting. The ReAct framework~\cite{ref28} combines reasoning and action within a single loop, allowing LLM agents to think through a problem while interacting with external tools. Reflexion~\cite{ref29} extends this idea with self-reflection, in which agents review their earlier decisions and use natural-language feedback to improve future performance. Toolformer~\cite{ref30} shows that language models can learn to use external tools autonomously, making them more capable on real-world tasks.

Chain-of-Thought prompting~\cite{ref31} encourages models to reason step-by-step through complex problems, significantly improving performance on tasks that require deeper inference. A broader survey of LLM-based autonomous agents~\cite{ref32} catalogues the common design elements---memory systems, strategy planning, and tool usage---across a large body of contemporary work. Signals builds directly on these ideas by combining the ReAct approach for agentic coordination with self-improvement concepts inspired by Reflexion, re-interpreted through a reinforcement-learning lens and implemented in the specific context of risk-aware memecoin trading.

\section{Sentiment Analysis for Crypto Markets}
Social-media sentiment has demonstrated genuine predictive power for cryptocurrency prices. Twitter sentiment achieves approximately 74\% directional accuracy in crypto price prediction~\cite{ref8}, and Kraaijeveld and De Smedt~\cite{ref9} establish Granger-causal relationships between Twitter sentiment and cryptocurrency returns for multiple large-cap tokens. FinBERT~\cite{ref40} provides a domain-adapted sentiment analysis model for financial text, while bi-LSTM architectures combined with RoBERTa~\cite{ref41} achieve strong performance on crypto-specific sentiment classification.

Multi-modal fusion approaches such as DASF-Net~\cite{ref42} and MSGCA~\cite{ref22} integrate text, price and volume Signals for improved forecasting. Research on pump-and-dump detection via social media~\cite{ref43} and Telegram NLP~\cite{ref44} further highlights the importance of social Signals, particularly in markets where manipulation is organised through private communication channels. However, sentiment systems in the literature typically operate independently of on-chain analysis, and very few attempt to incorporate sentiment into a real-time risk-adjusted trading decision.

Table~\ref{tab:lit} summarises the key findings of a subset of the most relevant prior work and positions Signals against the state of the art. Signals is, to the best of the authors' knowledge, the only system that simultaneously integrates market Signals, on-chain safety analysis, and social sentiment through a multi-agent architecture with self-improving memory, specifically targeting memecoin markets with real-time processing capabilities.

\begin{small}
\begin{longtable}{@{}p{0.8cm}p{3.8cm}p{3.4cm}p{5.6cm}@{}}
\caption{Literature Survey summary.}\label{tab:lit}\\
\toprule
\textbf{SR} & \textbf{Paper Title} & \textbf{Authors \& Venue} & \textbf{Key Findings} \\
\midrule
\endfirsthead
\multicolumn{4}{c}{\textit{(Table \ref{tab:lit} continued from previous page)}}\\
\toprule
\textbf{SR} & \textbf{Paper Title} & \textbf{Authors \& Venue} & \textbf{Key Findings} \\
\midrule
\endhead
\midrule
\multicolumn{4}{r}{\textit{(continued on next page)}}\\
\endfoot
\bottomrule
\endlastfoot
1. & Comprehensive review of deep learning for cryptocurrency price prediction & H. Wang \textit{et al.} (arXiv 2024)~\cite{ref7} & Surveys LSTM, transformer and hybrid architectures for crypto forecasting; identifies memecoins as an open problem. \\
2. & XGBoost-based cryptocurrency price prediction & arXiv 2024~\cite{ref4} & Gradient-boosted trees deliver competitive short-horizon forecasts but rely purely on technical indicators. \\
3. & Temporal Fusion Transformers for crypto prediction & arXiv 2025~\cite{ref21} & Multi-horizon attention-based forecaster fusing static and temporal features; ignores on-chain risk. \\
4. & MemeChain: multi-agent simulation of memecoin dynamics & arXiv 2026~\cite{ref1} & Simulation of memecoin market participants with emergent hype cycles; no real-time integration. \\
5. & ME2F: Memecoin ecosystem fragility analysis & arXiv 2025~\cite{ref2} & Identifies microstructural patterns that precede collapse events in memecoin markets. \\
6. & SolRPDS: Solana rug-pull detection system & arXiv 2025~\cite{ref3} & First dedicated Solana rug-pull dataset and detection pipeline. \\
7. & CRPWarner: cryptocurrency rug-pull detection with warning & arXiv 2024~\cite{ref5} & Achieves 91.8\% rug-pull detection accuracy using smart-contract features. \\
8. & RPHunter: behavioural rug-pull detection & arXiv 2025~\cite{ref12} & 95.3\% accuracy combining behavioural and contract-level features. \\
9. & TokenScout: early detection of DeFi token scams & ACM CCS 2024~\cite{ref13} & 98\% accuracy on live token behaviour evaluation. \\
10. & SoK: Understanding Rug Pulls in DeFi & arXiv 2024~\cite{ref20} & Catalogues 35 distinct rug-pull attack patterns; most tools detect only 25--63\%. \\
11. & TradingAgents: multi-agent LLM trading framework & Xiao \textit{et al.} (arXiv 2024)~\cite{ref23} & Debate-style multi-agent LLM framework for stock trading decisions. \\
12. & FinCon: synthesised multi-agent system with verbal reinforcement & Yu \textit{et al.} (arXiv 2024)~\cite{ref24} & Uses conceptual verbal reinforcement for portfolio management. \\
13. & CryptoTrade: reflective LLM crypto-trading agent & Li \textit{et al.} (arXiv 2024)~\cite{ref25} & Zero-shot reflective LLM crypto-trader, no on-chain safety checks. \\
14. & FinMem: layered-memory financial LLM agent & Yu \textit{et al.} (arXiv 2023)~\cite{ref26} & Memory-augmented LLM for financial decision making. \\
15. & ReAct: reasoning and acting in language models & Yao \textit{et al.} (ICLR 2023)~\cite{ref28} & Unified reasoning + acting loop enabling tool use in LLMs. \\
16. & Reflexion: verbal reinforcement learning agents & Shinn \textit{et al.} (NeurIPS 2023)~\cite{ref29} & Self-reflective natural-language feedback for continual improvement. \\
17. & Chain-of-Thought prompting & Wei \textit{et al.} (NeurIPS 2022)~\cite{ref31} & Step-by-step reasoning traces improve complex-problem performance. \\
18. & Survey on LLM-based autonomous agents & Wang \textit{et al.} (2024)~\cite{ref32} & Catalogues design dimensions: memory, planning, tool use. \\
19. & Twitter sentiment for crypto prediction & arXiv 2023~\cite{ref8} & 74\% directional accuracy using Twitter sentiment alone. \\
20. & Kraaijeveld \& De Smedt on Twitter sentiment & J. Int. Fin. Markets 2020~\cite{ref9} & Establishes Granger causality between sentiment and crypto returns. \\
21. & FinBERT: financial sentiment with pre-trained LMs & Araci (arXiv 2019)~\cite{ref40} & Domain-adapted BERT for financial sentiment classification. \\
22. & Pump.fun multimodal launch analysis & arXiv 2024~\cite{ref10} & Combines text, numeric and chart Signals for early-stage detection. \\
23. & Tulving's episodic-semantic memory model & 1972~\cite{ref45} & Cognitive-psychology foundation for Signals's dual memory. \\
24. & Pump-and-dump detection on social media & ACM TOIT 2023~\cite{ref43} & Detects coordinated manipulation via Twitter activity patterns. \\
25. & Telegram NLP pump-and-dump detection & IEEE BigData 2023~\cite{ref44} & Private-channel manipulation detection using NLP. \\
\end{longtable}
\end{small}

% =============================================================================
% CHAPTER 3 -- LIMITATIONS OF EXISTING SYSTEM
% =============================================================================
\chapter{Limitations of Existing System}
\label{ch:limit}

The literature survey in Chapter~\ref{ch:lit} reveals that despite substantial advances in cryptocurrency analytics, existing systems suffer from a set of consistent and compounding limitations when applied to the memecoin setting. These limitations are summarised below and collectively motivate the need for the Signals platform.

\begin{itemize}[leftmargin=8mm]
\item \textbf{Siloed analytical pipelines.} Price-prediction models, rug-pull detectors and sentiment analysers operate as independent tools. A trader who wishes to combine them is forced to juggle multiple dashboards, API keys and heterogeneous output formats, and must manually synthesise contradictory Signals under severe time pressure.

\item \textbf{Absence of memecoin-specific modelling.} Nearly all price-prediction research targets large-cap, liquid coins such as Bitcoin and Ethereum. The extraordinary volatility patterns of memecoins, their sensitivity to social hype cycles, and the prevalence of coordinated manipulation are fundamentally incompatible with assumptions baked into traditional time-series models.

\item \textbf{Rug detection treated as offline classification.} Existing rug-pull detectors~\cite{ref5,ref12,ref13} achieve high accuracy on retrospective datasets, but they are not integrated into live trading decision loops. A detector that flags a token as high-risk \emph{after} the liquidity has already been removed provides no protection to the retail investor.

\item \textbf{Sentiment analysis disconnected from on-chain truth.} Social-sentiment tools aggregate Twitter or Telegram mentions but rarely reconcile them with on-chain metrics such as holder concentration or liquidity lock status. A token that trends on Twitter with 90\% positive sentiment but has its liquidity unlocked should trigger an immediate warning---yet no single system in the surveyed literature makes this kind of cross-domain reconciliation automatic.

\item \textbf{Static heuristics with no self-improvement.} Most existing trading systems rely on hand-crafted rules or models that, once trained, do not adapt to changing market regimes. Memecoin markets transition rapidly between bull, bear, sideways and volatile states, rendering static policies brittle and short-lived.

\item \textbf{Opaque reinforcement-learning policies.} Deep-RL-based crypto trading agents~\cite{ref14,ref15,ref16} produce policies that are difficult to interpret or audit. For retail traders who need to understand \emph{why} a system is recommending a particular action, this opacity is a significant adoption barrier and a potential source of regulatory risk.

\item \textbf{Lack of real-time latency guarantees.} Memecoins can double or halve in price within minutes. Systems that require long retraining cycles, batch inference, or human-in-the-loop annotation cannot operate on the timescales required by live memecoin trading.

\item \textbf{Limited social-data diversity.} Even the strongest sentiment pipelines typically rely on a single data source, most often Twitter. Memecoin hype, however, also originates on Telegram, Discord, Reddit, TikTok and YouTube. Systems tied to one source miss Signals from the others.

\item \textbf{No standardised evaluation environment.} There is no widely accepted, reproducible simulator for memecoin trading. Each research group builds its own ad-hoc back-tester, which makes comparing methods across papers essentially impossible.
\end{itemize}

These limitations across current analytical tools highlight a critical research gap:
the complete absence of a unified, real-time, and self-improving platform
engineered specifically for the extreme volatility of modern memecoin markets.
Current systems lack multi-signal integration and transparent interpretability,
leaving traders heavily exposed to sudden and unpredictable momentum shifts.
The remainder of this report details exactly how the Signals architecture
systematically addresses each of these technical and operational gaps in turn.
The system is designed to remain highly resilient against market anomalies,
preventing technical mishaps during periods of extreme network congestion.
Signals translates chaotic market noise into actionable, interpretable data,
empowering traders and investors to operate with unprecedented efficiency.
This enables the memecoin community to confidently execute profitable trades
while rigorously mitigating their exposure to severe downside financial risks.

% =============================================================================
% CHAPTER 4 -- PROBLEM STATEMENT, OBJECTIVES AND SCOPE
% =============================================================================
\chapter{Problem Statement, Objectives and Scope}
\label{ch:prob}

\section{Problem Statement}
The memecoin ecosystem exhibits extreme price volatility, a high prevalence of rug-pull exploits, and sentiment-driven price dynamics that no existing analytical system unifies into a single real-time decision pipeline. As a direct result, retail investors are forced to rely on fragmented, siloed tools that individually address price prediction, contract safety, or social sentiment, but never all three together, leaving them exposed to fraud, manipulation, and avoidable losses. There is a pressing need for a unified, self-improving, multi-signal platform that can ingest live market data, on-chain safety metrics and social sentiment, reason about them through specialised agentic intelligence, and produce a risk-aware trading decision within a 30-second latency budget.

\section{Objectives}
The objectives of this project are as follows:

\begin{enumerate}[leftmargin=12mm]
\item To design and implement a multi-agent LLM architecture that simultaneously analyses market Signals, on-chain safety metrics and social sentiment for memecoin tokens.
\item To engineer a feature set of twenty-three domain-specific indicators spanning technical (11), on-chain (7) and social (5) dimensions, and to validate these features on live Solana tokens.
\item To develop a composite scoring framework with weighted signal fusion (40/40/20) and cross-domain safety overrides that prevent favourable market Signals from masking catastrophic on-chain risk.
\item To build a ReAct-based orchestration pipeline that autonomously queries external data APIs (DexScreener, GeckoTerminal, RugCheck, GoPlus, Twitter v2) under strict latency and resilience constraints.
\item To implement a reinforcement-learning agentic trader that executes a THINK-ACT-REFLECT-LEARN loop with dual memory (episodic trading journal + semantic rule base) and learns trading rules expressed in natural language without gradient-based training.
\item To develop a synthetic market simulator based on Geometric Brownian Motion with regime switching and jump diffusion, providing a controlled and reproducible evaluation environment.
\item To benchmark a representative set of open-source and proprietary LLMs on the same simulated trading task in order to quantify the effect of model choice on trading performance.
\item To ensure the entire system produces end-to-end trading decisions with latency under thirty seconds, making it practical for live memecoin markets.
\end{enumerate}

\section{Scope}
The scope of Signals covers the following functional areas:

\begin{itemize}[leftmargin=8mm]
\item \textbf{Supported blockchains and Network Architecture.} The system fundamentally targets high-throughput, low-latency networks—specifically Solana and Base—which currently serve as the primary epicenters for memecoin liquidity and hyper-volatile trading activity. This is achieved through a proprietary ChainCapabilities abstraction layer, which dynamically evaluates network conditions, standardizes disparate RPC node responses, and dictates the appropriate data sources and routing protocols at runtime without requiring core codebase alterations.

\item \textbf{Data sources.} The framework relies on a robust, multi-faceted data pipeline designed to capture the full spectrum of market mechanics. Quantitative market data, including tick-by-tick price action and liquidity depth, is continuously ingested from DexScreener and GeckoTerminal. Crucial on-chain safety metrics—such as contract renouncements, mint authorities, and honeypot risks—are audited in real-time via RugCheck and GoPlus. Simultaneously, qualitative social momentum and sentiment velocity are extracted via the Twitter/X v2 API. 

\item \textbf{Prediction horizon.} The system is engineered to produce immediate (minutes-to-hours) actionable buy, sell, or hold directives. These outputs are tightly coupled with granular confidence scores and categorical risk labels (e.g., Low, Medium, Extreme). The project explicitly excludes long-horizon price forecasting, as traditional long-term predictive models mathematically break down when applied to hype-driven decentralized assets.

\item \textbf{Trading modes.} The agentic trader currently operates exclusively within a secure, zero-risk simulation mode. In this environment, the ReAct-based agent executes and evaluates its strategies against a custom-built synthetic market simulator that mimics slippage, network congestion, and volatile price swings. Fully automated, live trading directly interfacing with decentralized exchange infrastructures (specifically utilizing the Jupiter and Raydium SDKs for Solana) requires advanced transaction signing and slippage tolerance management; consequently, this is architected strictly as a deployment extension and is detailed comprehensively in the future-work chapter.

\item \textbf{User deliverables.} The user receives a structured JSON decision object containing the action, confidence, risk level, reasoning trace, and any triggered learned rules, along with a dashboard view summarising the contributing Signals.

\item \textbf{Ethical use.} Signals is designed as a decision-support tool for educational and research purposes. It is not an automated trading bot intended for mass deployment without explicit user consent and risk acknowledgement.

\item \textbf{Evaluation scope.} Evaluation uses both the synthetic market simulator and a live Solana token case study (MURICA) to validate end-to-end latency and decision quality under real-world conditions.

\item \textbf{Sustainability alignment.} Signals's contribution to decent work (SDG~8), industry innovation (SDG~9), and peace and strong institutions through fraud prevention (SDG~16) is discussed in Chapter~\ref{ch:impl}.
\end{itemize}

% =============================================================================
% CHAPTER 5 -- PROPOSED SYSTEM
% =============================================================================
\chapter{Proposed System}
\label{ch:proposed}

The proposed Signals system is a decentralised multi-agent architecture designed to unify memecoin intelligence across market, on-chain and social domains. At a formal level, the complete system is defined as the triple $\mathcal{S} = \langle \mathcal{D}, \mathcal{A}, \mathcal{T} \rangle$, where $\mathcal{D}$ denotes the Data Acquisition Layer, $\mathcal{A}$ denotes the Multi-Agent Analysis Layer, and $\mathcal{T}$ denotes the Trading Layer. Each layer is designed to be independently replaceable, horizontally scalable, and resilient to failures in any of its sub-components.

\section{Modules}
Signals comprises the following core modules:

\begin{enumerate}[leftmargin=12mm]
\item \textbf{Data Acquisition Module.} Connects to five external APIs in parallel via Python's \texttt{asyncio.gather}, implementing circuit breakers, exponential backoff and fallback strategies.
\item \textbf{Feature Engineering Module.} Transforms raw API responses into twenty-three normalised features across technical, on-chain and social domains.
\item \textbf{Multi-Agent Analysis Module.} Runs four specialised LLM agents (Market, Safety, Social, Prediction) with distinct temperature profiles and structured JSON output parsing.
\item \textbf{Composite Scoring Engine.} Fuses agent outputs using a weighted formula ($40\%$ technical + $40\%$ on-chain + $20\%$ social) with cross-domain safety overrides.
\item \textbf{Agentic Trader Module.} Executes the THINK-ACT-REFLECT-LEARN cycle using a dual memory architecture and autonomously discovers natural-language trading rules.
\item \textbf{Synthetic Market Simulator.} Provides controlled evaluation through Geometric Brownian Motion with regime switching and jump diffusion.
\item \textbf{Resilience Manager.} Implements three-tier fault tolerance at the API, retry and orchestration levels, including graceful confidence degradation when data sources are partially unavailable.
\item \textbf{User Dashboard.} Presents the final trading decision, risk level, confidence and contributing Signals in a human-readable format.
\end{enumerate}

\section{System Architecture}
Figure~\ref{fig:arch} shows the overall Signals architecture. A token address enters the five-phase ReAct orchestrator, which fetches data from four external API sources in parallel with circuit-breaker protection. Twenty-three features are engineered and analysed by four specialised LLM agents with distinct temperature profiles. The Prediction Agent synthesises a composite score with cross-domain safety overrides, which feeds the agentic trader's THINK-ACT-REFLECT-LEARN loop with dual memory.

\begin{figure}[H]
    \centering
    \includegraphics[width=\textwidth]{architecture.jpg}
    \caption{Data Flow Diagram -- Level 0}
    \label{fig:dfd0}
\end{figure}



\section{UML Diagrams}

\subsection{Data Flow Diagram -- Level 0}
Figure~\ref{fig:dfd0} shows the Level 0 Data Flow Diagram of Signals. At this highest level of abstraction the system has three external entities: the \emph{User}, the \emph{Signals System}, and the \emph{External Data Providers}. The user submits a token address; the system queries external providers, processes the responses, and returns a decision object.

\begin{figure}[H]
    \centering
    \includegraphics[width=\textwidth]{dfdlevel0.png}
    \caption{Data Flow Diagram -- Level 0}
    \label{fig:dfd0}
\end{figure}

\subsection{Data Flow Diagram -- Level 1}
Figure~\ref{fig:dfd1} decomposes the Signals System into its three internal layers: Data Acquisition, Multi-Agent Analysis, and Trading. Arrows show how feature vectors flow from acquisition to analysis, how agent outputs flow to the composite scoring engine, and how trading decisions flow back to the user.

At the Data Acquisition layer, heterogeneous inputs—including market prices, liquidity metrics, on-chain Signals, and social sentiment—are collected, normalised, and transformed into structured feature vectors. This layer ensures consistency in data formats and timestamps before forwarding them downstream.

Overall, the Level 1 DFD highlights a clear separation of concerns and a unidirectional yet feedback-aware data flow, ensuring modularity, interpretability, and scalability of the system.
\begin{figure}[H]
    \centering
    \includegraphics[width=\textwidth]{dfdlevel1.png}
    \caption{Data Flow Diagram -- Level 1}
    \label{fig:dfd1}
\end{figure}

\subsection{Data Flow Diagram -- Level 2}
Figure~\ref{fig:dfd2} provides the Level 2 decomposition, showing the individual API connectors within the Data Layer, the per-agent prompt templates within the Analysis Layer, and the episodic/semantic memory stores within the Trading Layer. Each block has its own data store for caching, rate-limit tracking, and historical logging.

At the Data Layer, connectors for sources such as market APIs, on-chain analytics, and social platforms operate independently, enabling parallel data ingestion and fault isolation. Each connector maintains its own cache and retry logic, ensuring robustness against API failures and rate limits.  

Within the Analysis Layer, agent-specific prompt templates define structured reasoning for Market, Safety, Social, and Prediction agents. These agents process normalized inputs and generate validated outputs, allowing consistent communication across modules.  

Overall, the Level 2 DFD highlights a modular and decoupled architecture, where independent components interact through well-defined interfaces, improving scalability, maintainability, and fault tolerance.

\begin{figure}[H]
    \centering
    \includegraphics[width=125mm]{dfdlevel2.png}
    \caption{Data Flow Diagram -- Level 2}
    \label{fig:dfd2}
\end{figure}
\subsection{Class Diagram}
The class diagram (Figure~\ref{fig:class}) shows the principal Python classes: \texttt{SignalsOrchestrator}, \texttt{DataAcquisition}, \texttt{FeatureEngineer}, the four \texttt{Agent} subclasses, \texttt{CompositeScorer}, \texttt{AgenticTrader}, \texttt{EpisodicMemory}, \texttt{SemanticRuleStore}, \texttt{WalletManager}, and the \texttt{SyntheticMarketSimulator}. Inheritance relationships, association multiplicities, and key method signatures are annotated.

\begin{figure}[H]
    \centering
    \includegraphics[width=0.68\textwidth]{class.png}
    \caption{Signals class diagram}
    \label{fig:class}
\end{figure}

\subsection{Sequence Diagram}
Figure~\ref{fig:seq} shows the sequence of messages exchanged between the User, the Orchestrator, the Data Acquisition Layer, the four Agents, and the Agentic Trader when a new token address is submitted. The user request triggers parallel API calls, followed by sequential agent invocation, composite scoring, and finally the agentic trading decision.

\begin{figure}[H]
    \centering
    \includegraphics[width=\textwidth]{sequence.jpg}
    \caption{Sequence diagram for a single token analysis request}
    \label{fig:seq}
\end{figure}

\subsection{Activity Diagram}
The activity diagram (Figure~\ref{fig:activity}) visualises the overall workflow of a Signals analysis session, from user input through data acquisition, per-agent analysis, cross-domain safety overrides, composite scoring, and agentic trading to final decision delivery.

\begin{figure}[H]
    \centering
    \includegraphics[width=0.6\textwidth]{activity.png}
    \caption{Activity diagram of the Signals pipeline}
    \label{fig:activity}
\end{figure}

\section{Methodology}

\subsection{Data Acquisition Layer}
The data acquisition layer attaches to five external APIs grouped into three signal categories (Table~\ref{tab:api}). Market data such as price, volume, liquidity and trading pairs is fetched from DexScreener and GeckoTerminal. Safety data such as rug score, honeypot flags and holder distribution is fetched from RugCheck (Solana) and GoPlus (multi-chain). Social data consisting of mentions, sentiment scores and influencer activity is fetched from the Twitter v2 API. Each API connection is safeguarded by a circuit breaker with three states (\textsc{closed}, \textsc{open}, \textsc{half-open}); after three consecutive failures the circuit opens and returns a fallback value for 30 seconds before attempting recovery. Data is fetched in parallel using Python's \texttt{asyncio.gather} with a 30-second timeout.

\begin{table}[H]
\centering
\caption{API Data Sources and Configuration.}
\label{tab:api}
\begin{tabular}{@{}llll@{}}
\toprule
\textbf{API} & \textbf{Type} & \textbf{Key Fields} & \textbf{Timeout} \\
\midrule
DexScreener & Market & Price, Volume, Pairs & 10 s \\
GeckoTerminal & Market & Market Cap, Liquidity & 10 s \\
RugCheck & Safety & Rug Score (Solana) & 30 s \\
GoPlus & Safety & Honeypot, Holders & 30 s \\
Twitter v2 & Social & Mentions, Sentiment & 30 s \\
\bottomrule
\end{tabular}
\end{table}

\subsection{Multi-Agent Analysis Layer}
The analysis layer contains four specialised LLM agents, each executing with a principled temperature profile (Table~\ref{tab:agents}). Lower temperatures (0.1--0.2) are used for safety-critical and final prediction tasks to ensure conservative, deterministic outputs, while higher temperatures (0.3--0.4) are used for market and social analysis to permit creative pattern recognition in ambiguous inputs. Each agent emits structured JSON validated via Pydantic models, with a fallback path that attempts to extract JSON objects from markdown fences or bare string responses. A retry mechanism with error-feedback prompts recovers from malformed LLM outputs.

\begin{table}[H]
\centering
\caption{Multi-Agent LLM Configuration.}
\label{tab:agents}
\begin{tabular}{@{}llcp{6cm}@{}}
\toprule
\textbf{Agent} & \textbf{Model} & \textbf{Temp} & \textbf{Key Output} \\
\midrule
Market     & Gemini 2.0 Flash & 0.3 & Market health (1--10), volume \& momentum score \\
Safety     & Llama 3          & 0.2 & Rug risk (0--100), risk flags, SAFE/CAUTION/AVOID \\
Social     & Grok 4           & 0.4 & Sentiment (0--100), hype level, manipulation flag \\
Prediction & DeepSeek R3      & 0.1 & BUY/SELL/HOLD decision, confidence, rationale \\
\bottomrule
\end{tabular}
\end{table}

\subsection{Feature Engineering}
The system engineers twenty-three features across three categories, listed in Table~\ref{tab:feat}.

\begin{table}[H]
\centering
\caption{Complete 23-Feature Catalogue.}
\label{tab:feat}
\small
\begin{tabular}{@{}clll@{}}
\toprule
\textbf{\#} & \textbf{Feature} & \textbf{Category} & \textbf{Range} \\
\midrule
1  & price                      & Technical & $\mathbb{R}_+$ \\
2  & price change 1h            & Technical & $[-100, \infty)$\,\% \\
3  & price change 24h           & Technical & $[-100, \infty)$\,\% \\
4  & volume 24h                 & Technical & $\mathbb{R}_+$ \\
5  & market cap                 & Technical & $\mathbb{R}_+$ \\
6  & liquidity                  & Technical & $\mathbb{R}_+$ \\
7  & RSI(14)                    & Technical & $[0, 100]$ \\
8  & MACD(12,26,9)              & Technical & $\mathbb{R}$ \\
9  & Bollinger(20,2)            & Technical & $[-1, 1]$ \\
10 & volatility                 & Technical & $[0, 1]$ \\
11 & volume ratio               & Technical & $\mathbb{R}_+$ \\
12 & holder count               & On-Chain  & $\mathbb{Z}_+$ \\
13 & top-10 holder pct          & On-Chain  & $[0, 100]$\,\% \\
14 & smart money flow           & On-Chain  & $\{\mathrm{buy}, \mathrm{sell}, \mathrm{neutral}\}$ \\
15 & rug score                  & On-Chain  & $[0, 100]$ \\
16 & dev wallet pct             & On-Chain  & $[0, 100]$\,\% \\
17 & liquidity locked           & On-Chain  & $\{0, 1\}$ \\
18 & lock days remaining        & On-Chain  & $\mathbb{Z}_{\geq 0}$ \\
19 & mentions 24h               & Social    & $\mathbb{Z}_{\geq 0}$ \\
20 & sentiment score            & Social    & $[0, 100]$ \\
21 & influencer mentions        & Social    & $\mathbb{Z}_{\geq 0}$ \\
22 & trending                   & Social    & $\{0, 1\}$ \\
23 & community size             & Social    & $\mathbb{Z}_+$ \\
\bottomrule
\end{tabular}
\end{table}

\subsection{ReAct Orchestration}
The orchestrator implements a five-phase ReAct pipeline that coordinates data acquisition, validation, analysis and synthesis. In the \emph{Think} phase the orchestrator plans which APIs to call given the target chain. In the \emph{Act} phase it issues parallel API calls. In the \emph{Observe} phase it validates the responses against expected schemas. In the second \emph{Act} phase the specialised agents perform their analyses. Finally, in the \emph{Reflect} phase the prediction agent synthesises the decision. This explicit phase structure provides transparent execution traces that can be logged, audited, and replayed for debugging.

\subsection{Composite Scoring Framework}
The prediction pipeline employs a weighted composite scoring framework that fuses outputs from all three analytical domains. The per-domain scores are computed as:

\begin{align}
S_{\mathrm{tech}}   &= \frac{f_{\mathrm{RSI}} + f_{\mathrm{MACD}} + f_{\mathrm{BB}}}{3}\\
S_{\mathrm{chain}}  &= 0.5 \cdot f_{\mathrm{rug}} + 0.3 \cdot f_{\mathrm{liq}} + 0.2 \cdot f_{\mathrm{smart}}\\
S_{\mathrm{social}} &= 0.7 \cdot f_{\mathrm{sent}} + 0.3 \cdot f_{\mathrm{trend}}\\
S_{\mathrm{overall}} &= 0.4 \cdot S_{\mathrm{tech}} + 0.4 \cdot S_{\mathrm{chain}} + 0.2 \cdot S_{\mathrm{social}}
\end{align}

Cross-domain safety overrides are formalised as a piecewise confidence modifier:
\begin{equation}
c_{\mathrm{final}} = \begin{cases}
0.1 & \text{if honeypot detected}\\
c_{\mathrm{adj}} \cdot 0.5 & \text{if rug risk} > \theta_{\mathrm{rug}}\\
1.0 & \text{otherwise}
\end{cases}
\end{equation}
where $\theta_{\mathrm{rug}} = 70$ is the rug-risk threshold. This ensures that high-risk Signals dominate regardless of otherwise favourable market or social indicators. The on-chain risk score itself is computed from multiple binary and continuous indicators:
\begin{equation}
R_{\mathrm{total}} = \min\left(\sum_{i=1}^{n} w_i \cdot \mathbb{1}[\mathrm{risk}_i],\ 100\right)
\end{equation}
where the risk factors and their weights are: honeypot ($w = 40$), mintable token ($w = 15$), ownership not renounced ($w = 10$), top-10 holder concentration $>50\%$ ($w = 20$) or $>30\%$ ($w = 10$), and liquidity not locked ($w = 15$). Risk levels are mapped as: $0$--$20$ = LOW, $21$--$45$ = MEDIUM, $46$--$70$ = HIGH, $71$--$100$ = CRITICAL.

\subsection{Agentic Trading Loop}
The agentic trader begins each cycle by executing the THINK-ACT-REFLECT-LEARN loop depicted in Figure~\ref{fig:agentloop}.

\begin{figure}[H]
    \centering
    \includegraphics[width=\textwidth]{cycle.jpg}
    \caption{Agentic Trading loop}
    \label{fig:agentloop}
\end{figure}


\textbf{THINK.} A TokenSnapshot containing all twenty-three features, the current wallet position, the last $n$ trades from episodic memory, and the currently active rules from semantic memory is passed to an LLM agent (Gemini 2.0 Flash, temperature~0.3). The agent produces a Trade Decision comprising an action (\textsc{buy}, \textsc{sell} or \textsc{hold}), a confidence score in $[0, 100]$, a natural-language rationale, a risk assessment, and optional target and stop-loss prices.

\textbf{ACT.} Execution is performed via a WalletManager subject to three constraints: (1) per-trade position size is capped at 25\% of the total portfolio balance ($p_{\max} = 0.25$); (2) an automatic stop-loss is placed at 10\% below the weighted average entry price; and (3) the weighted average entry price is updated across multi-buy events for correct P\&L accounting.

\textbf{REFLECT.} After every trade the LLM performs a post-trade analysis, comparing the original rationale against the observed outcome. The outcome is classified as \textsc{success}, \textsc{failure} or \textsc{neutral}, and key learning points are extracted.

\textbf{LEARN.} Reflections are distilled into long-term TradingRules belonging to one of three categories: \textsc{entry} (conditions under which to buy), \textsc{exit} (conditions under which to sell) and \textsc{avoid} (conditions that should block any entry). Each rule is stored in Markdown with an ID, category, title, description, activation conditions, success and failure counts, cumulative P\&L, and the IDs of the trades that generated it. Rules are loaded back into future THINK prompts via retrieval-augmented contexts.

\subsection{Dual Memory Architecture}
Persistent state is partitioned into two independent repositories modelled on Tulving's distinction between episodic and semantic memory~\cite{ref45}.

\paragraph{Episodic Memory.} A chronological log of \texttt{TradeRecord} instances, each with fourteen attributes: trade identifier, timestamp, token name and symbol, action, entry and exit price, quantity, dollars traded, realised profit or loss (in dollars and percentage), original rationale, risk assessment, outcome label and the rule IDs that influenced the trade. The most recent $n$ records are injected into the explicit prompt context in the THINK phase.

\paragraph{Semantic Memory.} A collection of \texttt{TradingRule} objects of type \textsc{entry}, \textsc{exit} or \textsc{avoid}. Each rule tracks its success count $s$, failure count $f$, cumulative profit and loss, and trade identifiers on which it was founded. The win rate $\pi = s / (s + f)$ serves as an evidence-based weight used for rule retrieval, conflict resolution and pruning.

\subsection{Synthetic Market Simulator}
To enable controlled experimentation, a synthetic market simulator is developed based on Geometric Brownian Motion with regime switching and jump diffusion. The price dynamics follow:
\begin{equation}
\frac{dP}{P} = \mu_r\, dt + \sigma_r\, dW + J\, dN(\lambda)
\end{equation}
where $\mu_r$ and $\sigma_r$ are the drift and volatility for regime $r$, $W$ is a standard Wiener process, $J \sim \mathcal{U}(-0.15, 0.20)$ is the jump magnitude, and $N(\lambda)$ is a Poisson process with intensity $\lambda = 0.05$ (i.e.\ 5\% probability per step).

The simulator implements Markov regime switching with four states (Table~\ref{tab:regimes}). Regimes transition with weights $w = [0.25,\,0.25,\,0.35,\,0.15]$, and regime durations follow a uniform distribution $\mathcal{U}(50, 150)$ steps. This generates realistic market dynamics including sustained trends, consolidation periods and unexpected volatility spikes. On-chain metrics are generated with configurable correlation to price movement: holder count evolves as $h_t = h_{t-1} + \mathcal{U}(-10, 5)$ in bear regimes and $h_{t-1} + \mathcal{U}(-5, 10)$ in bull regimes. Top-10 concentration is modelled as $c_{10} = 30 + \mathcal{U}(-5, 5) + (50 - h/100)$, clipped to $[20, 80]\%$. The rug score is drawn from a base range $[5, 50]$ with additional penalties: $+10$ if $c_{10} > 60\%$, $+15$ if liquidity is not locked, and $+10$ if the developer wallet exceeds 10\%. Social metrics correlate with price dynamics: sentiment $= 50 + p_{24h} \cdot 2 + \mathcal{U}(-10, 10)$, clipped to $[0, 100]$; mentions scale with community size and price changes; and influencer mentions follow a Poisson distribution conditioned on sentiment. Finally, the technical indicators RSI-14, MACD-12/26/9 and Bollinger Bands-20/2 are computed from the simulated price history rather than directly approximated, ensuring realistic indicator behaviour.

\begin{table}[H]
\centering
\caption{Synthetic Market Regime Parameters.}
\label{tab:regimes}
\begin{tabular}{@{}lccc@{}}
\toprule
\textbf{Regime} & \textbf{Drift ($\mu$)} & \textbf{Volatility ($\sigma$)} & \textbf{Weight} \\
\midrule
Bull      & $+0.002$  & $0.030$ & $0.25$ \\
Bear      & $-0.002$  & $0.030$ & $0.25$ \\
Sideways  & $\phantom{+}0.000$   & $0.015$ & $0.35$ \\
Volatile  & $\phantom{+}0.000$   & $0.080$ & $0.15$ \\
\bottomrule
\end{tabular}
\end{table}
\subsection{Resilience and Fault Tolerance}

Signals adopts a multi-layered resilience architecture designed to ensure stability under unreliable data sources, API failures, and adversarial conditions. The system is structured across three complementary tiers: API-level protection, retry-level recovery, and orchestration-level adaptation.

At the retry level, transient failures such as network timeouts or rate limits are handled through exponential backoff. The delay for the $k$-th retry is given by:
\begin{equation}
\tau_k = \min(\tau_0 \cdot b^{k-1}, \tau_{\max})
\end{equation}
where $\tau_0 = 1.0$\,s is the initial delay, $b = 2.0$ is the backoff factor, and $\tau_{\max} = 10.0$\,s is the maximum cap. This strategy reduces load on unstable services while increasing the probability of successful recovery without overwhelming external APIs.

At the orchestration level, Signals incorporates graceful degradation to maintain functionality under partial system failure. Instead of halting execution when some data sources are unavailable, the system dynamically adjusts prediction confidence based on the proportion of valid inputs:
\begin{equation}
c_{\mathrm{adj}} = \frac{|\text{valid sources}|}{|\text{total sources}|}
\end{equation}


Finally, the system logs all failure events, retry attempts, and circuit breaker transitions, enabling post-hoc analysis and continuous improvement of fault-handling strategies. Together, these mechanisms ensure that Signals remains reliable, adaptive, and robust in highly volatile and failure-prone crypto market environments.

% =============================================================================
% CHAPTER 6 -- EXPERIMENTAL SETUP
% =============================================================================
\chapter{Experimental Setup}
\label{ch:setup}

\section{Technology Stack}

\subsection{Backend}
\begin{itemize}[leftmargin=8mm]
\item \textbf{Language:} Python 3.11
\item \textbf{Async runtime:} \texttt{asyncio} with \texttt{aiohttp} for parallel API access
\item \textbf{LLM SDKs:} OpenAI, Google Gemini, Anthropic, xAI Grok, DeepSeek, Together AI
\item \textbf{Data validation:} Pydantic v2 for structured JSON model validation
\item \textbf{Technical indicators:} \texttt{pandas-ta} and a custom indicator module
\item \textbf{Storage:} SQLite for episodic memory, Markdown files for semantic rules
\item \textbf{Testing:} \texttt{pytest} with fixture-driven simulator tests
\end{itemize}

\subsection{Frontend / Dashboard}
\begin{itemize}[leftmargin=8mm]
\item \textbf{Framework:} Next.js 14 with React 18 and TypeScript
\item \textbf{Styling:} Tailwind CSS with shadcn/ui components
\item \textbf{Charts:} Recharts for candlestick and equity-curve visualisations
\item \textbf{State management:} TanStack Query for server state
\end{itemize}

\subsection{External Services}
\begin{itemize}[leftmargin=8mm]
\item DexScreener Public API
\item GeckoTerminal API
\item RugCheck.xyz API (Solana rug-score service)
\item GoPlus Security API (multi-chain honeypot and contract risk)
\item Twitter / X v2 API with elevated access for historical search
\end{itemize}

\section{Dataset and Inputs}
Experiments use two complementary data sources:

\begin{enumerate}[leftmargin=12mm]
\item \textbf{Synthetic Market Simulator.} A fully controllable environment generated by the regime-switching GBM described in Chapter~\ref{ch:proposed}. Each episode spans 50--500 steps, and all technical, on-chain and social metrics are generated consistently with the simulated price path.
\item \textbf{Live Solana Tokens.} A live case study uses the MURICA token (address \texttt{HQ4...murica}), which exhibited \$10.3M 24-hour volume and a 3{,}576\% 24-hour price change. Data is fetched from the production endpoints of DexScreener, RugCheck and Twitter v2.
\end{enumerate}

\section{Performance Evaluation Parameters}
The following metrics are used to evaluate Signals:

\begin{itemize}[leftmargin=8mm]
\item \textbf{Cumulative return (\%)} --- total portfolio growth over an episode.
\item \textbf{Win rate (\%)} --- fraction of closed trades with positive P\&L.
\item \textbf{Sharpe ratio} --- risk-adjusted return of the equity curve.
\item \textbf{Maximum drawdown} --- peak-to-trough decline in portfolio value.
\item \textbf{Number of trades} --- total trades executed per episode.
\item \textbf{Autonomous rules discovered} --- number of unique TradingRules learned.
\item \textbf{End-to-end latency} --- wall-clock time from token address to decision object.
\item \textbf{API success rate} --- fraction of external calls returning valid data within timeout.
\end{itemize}

\section{Software Requirements}
\begin{itemize}[leftmargin=8mm]
\item Python 3.11+, Node.js 20+, npm / pnpm
\item Git for version control; GitHub for remote
\item Visual Studio Code with Python, Pylance and Tailwind IntelliSense extensions
\item Docker (optional, for deployment)
\item API keys: OpenAI, Google, Anthropic, xAI, DeepSeek, Twitter v2
\end{itemize}

\section{Hardware Requirements}

To ensure efficient development, experimentation, and deployment of the proposed system, a balanced combination of local computing resources and cloud infrastructure is required. The following specifications outline a practical setup that supports both model development and real-time system execution.

\begin{itemize}[leftmargin=8mm]

\item \textbf{Development Laptop}
A modern development machine is essential for coding, data preprocessing, lightweight model experimentation, and running local services such as APIs and dashboards.

\begin{itemize}[leftmargin=8mm]
\item \textbf{Processor:} Intel Core i5 12th Gen or equivalent (6 cores, 12 threads)
This ensures smooth multitasking, faster data processing, and efficient handling of concurrent services such as backend servers and data pipelines.

\item \textbf{RAM:} 16 GB DDR4
Adequate memory is required for running multiple development tools simultaneously (e.g., IDEs, databases, containers, and browsers). It also supports moderate in-memory data processing and experimentation with machine learning models.

\item \textbf{Storage:} 512 GB SSD
Fast SSD storage significantly improves read/write speeds for datasets, model checkpoints, and containerized environments. It also reduces latency in development workflows.

\item \textbf{GPU:} NVIDIA RTX 3050 or better (optional)
While not mandatory, a dedicated GPU accelerates local experimentation with small language models (SLMs), deep learning architectures, and embedding generation. This is particularly useful for prototyping before scaling to cloud environments.

\item \textbf{Operating System:} Windows 11 / Ubuntu 22.04 / macOS 14
The system should support containerization (e.g., Docker), modern development frameworks, and compatibility with machine learning libraries.

\end{itemize}
\item \textbf{Network Requirements}
A stable, high-speed internet connection is essential for real-time data access and API communication.

\begin{itemize}[leftmargin=8mm]
\item \textbf{Latency:} $<200$ ms to US/EU endpoints for timely data aggregation.
\item \textbf{Bandwidth:} 50--100 Mbps recommended for smooth data flow and low-latency responses.
\end{itemize}

\item \textbf{Deployment Infrastructure}
Cloud infrastructure ensures scalability and reliability.

\begin{itemize}[leftmargin=8mm]
\item \textbf{Minimum Setup:} 4 vCPU, 8 GB RAM for basic testing and inference.
\item \textbf{Scalability:} Horizontal scaling with load balancing and orchestration (e.g., Kubernetes).
\item \textbf{Enhancements:} GPU instances for faster inference; caching (e.g., Redis) for improved performance.
\end{itemize}

Overall, this setup balances cost, performance, and scalability for real-time deployment.
% =============================================================================
% CHAPTER 7 -- IMPLEMENTATION AND RESULTS
% =============================================================================
\chapter{Implementation and Results}
\label{ch:impl}

\section{Implementation Overview}
The Signals reference implementation is written in Python~3.11 on the backend, with a Next.js~14 frontend for the user dashboard. The backend exposes a REST API (\texttt{POST /analyze}) that accepts a token address and an optional chain identifier, and returns a JSON decision object containing the action, confidence, risk level, agent rationales, and triggered rules. The implementation follows the module decomposition described in Chapter~\ref{ch:proposed}; the following sub-sections highlight the most important engineering details.

\subsection{Data Acquisition Module}
Each of the five external APIs is wrapped in a dedicated Python connector class that inherits from an abstract \texttt{BaseAPIConnector}. The base class implements circuit-breaker logic using the three-state model (\textsc{closed}, \textsc{open}, \textsc{half-open}), exponential backoff with jitter, and response caching keyed by token address and timestamp bucket. Parallel fetching is achieved through \texttt{asyncio.gather} with a 30-second overall timeout and per-connector timeouts as listed in Table~\ref{tab:api}.

\subsection{Feature Engineering Module}
Raw responses are normalised into a single \texttt{TokenSnapshot} dataclass containing all twenty-three features. Technical indicators (RSI-14, MACD-12/26/9, Bollinger Bands-20/2) are computed from the recent price history returned by DexScreener. On-chain features are extracted from RugCheck and GoPlus responses with defensive null-handling. Social features are derived from the Twitter v2 search API using query templates that match the token's ticker and contract address.

\subsection{Multi-Agent Module}
Each of the four agents is a subclass of a generic \texttt{Agent} class parameterised by model name, temperature, system prompt template and Pydantic output schema. The base class handles prompt rendering, LLM invocation, JSON extraction with fallback to markdown-fence parsing, Pydantic validation, and retry with error-feedback messages on validation failure.

\subsection{Agentic Trader Module}
The agentic trader is implemented as a single \texttt{AgenticTrader} class that owns references to an \texttt{EpisodicMemory} (SQLite-backed) and a \texttt{SemanticRuleStore} (Markdown-backed). On each THINK invocation, the most recent $n = 20$ trade records and all active rules are injected into the prompt context. Trade execution is delegated to a \texttt{WalletManager} that enforces the 25\% position-size cap and 10\% stop-loss.

\subsection{Synthetic Market Simulator}
The simulator is implemented as a generator that yields \texttt{TokenSnapshot} objects one step at a time. Regime state is maintained explicitly and transitions according to the weight vector described in Table~\ref{tab:regimes}. Jump events are sampled from a Poisson process with $\lambda = 0.05$.

\section{Implementation Screenshots}
The following placeholder figures illustrate representative screens from the implemented system. These are to be replaced with real screenshots during the final hardcopy submission.

\begin{figure}[H]
    \centering
    \includegraphics[width=\textwidth]{dashboard.jpg}
    \caption{Signals landing page}
    \label{fig:ss1}
\end{figure}


\begin{figure}[H]
    \centering
    \includegraphics[width=\textwidth]{Signal Scanner.jpeg}
    \caption{Autonomous Discovery Sweep with Kill-chain filter}
    \label{fig:ss2}
\end{figure}


\begin{figure}[H]
    \centering
    \includegraphics[width=\textwidth]{Decision.jpeg}
    \caption{Per-token analysis card showing Market, Rug Check, Social, Predictor agent scores and gives final action verdict}
    \label{fig:ss3}
\end{figure}


\begin{figure}[H]
    \centering
    \includegraphics[width=\textwidth]{agentictrader.jpg}
    \caption{Agentic trader dashboard showing episodic memory and learned rules.}
    \label{fig:ss4}
\end{figure}




\section{Experimental Results}

\subsection{Experimental Setup Recap}
Experiments use the synthetic market simulator of Section~5.4.8 with Gemini~2.0 Flash as the reasoning backbone for the agentic trader (temperature~0.3). Each agentic-trading episode is initialised with a balance of \$100; the multi-model benchmark uses \$10{,}000 per episode. Evaluation metrics include cumulative return (\%), win rate (\%), Sharpe ratio, maximum drawdown, number of trades and number of autonomous rules discovered. The Q-learning baseline uses a discretised state space over 10 episodes of 499 steps each. LLM API calls are made to the production endpoints of each provider; the estimated cost per 100-step agentic episode is roughly \$0.50--\$2.00 depending on the model chosen.

\subsection{Agentic Trader vs.\ Q-Learning Baseline}
Table~\ref{tab:agentvql} compares the agentic trader against a Q-learning baseline on the synthetic market simulator. The Q-learning baseline achieves only $-0.02\%$ average return, while the agentic trader attains $+9.86\%$ over medium-length episodes. Three factors explain this gap. First, the Q-learning agent must discretise the 23-dimensional feature vector into tabular states, losing critical information about feature interactions and continuous relationships, whereas the LLM agent processes the full feature vector in natural language. Second, Q-learning requires thousands of state-action visits to converge; the LLM leverages pre-trained financial knowledge and produces sensible decisions from the very first episode. Third, the LLM brings implicit priors about momentum, mean reversion and risk management that tabular RL lacks entirely.

\begin{table}[H]
\centering
\caption{Agentic Trader vs.\ Q-Learning Baseline.}
\label{tab:agentvql}
\begin{tabular}{@{}lccccc@{}}
\toprule
\textbf{Configuration} & \textbf{Ep.} & \textbf{Steps} & \textbf{Avg Ret.} & \textbf{Best} & \textbf{Trades} \\
\midrule
Agentic (medium)  & 2  & 100 & $+9.86\%$  & $+22.60\%$ & 5 \\
Agentic (short)   & 2  & 50  & $+0.80\%$  & $+5.20\%$  & 3 \\
Q-Learning        & 10 & 499 & $-0.02\%$  & $+0.31\%$  & 16 \\
\bottomrule
\end{tabular}
\end{table}

The agentic trader executes fewer trades (5 vs.\ 16) with greater conviction, demonstrating the kind of selective entry behaviour associated with professional traders. This conservative approach---with approximately 96\% HOLD decisions---reflects an unstated understanding that in highly volatile markets the cost of a wrong entry often exceeds the opportunity cost of inaction.

\begin{figure}[H]
    \centering
    \includegraphics[width=\textwidth]{fig4_equity_curves.png}
    \caption{Equity curves comparing the agentic trader against Q-learning.}
    \label{fig:eq}
\end{figure}


\subsection{Multi-Model Benchmark}
Eight LLMs were benchmarked as the reasoning backbone of the agentic trader (Table~\ref{tab:bench}). Several key findings emerge.

\begin{table}[H]
\centering
\caption{Open-Source vs.\ Proprietary LLM Trading Performance.}
\label{tab:bench}
\small
\begin{tabular}{@{}lrrrrrrrr@{}}
\toprule
\textbf{\#} & \textbf{Model} & \textbf{Return \%} & \textbf{Win \%} & \textbf{Sharpe} & \textbf{Trades} & \textbf{Max Dip} & \textbf{Long \%} & \textbf{Avg Hold} \\
\midrule
1 & Grok-4.20          & $+34.59$  & 31.6 & 0.019  & 158 & $-18.7\%$ & 85\% & 12.3 \\
2 & GPT-5.1            & $+8.88$   & 34.1 & 0.009  & 537 & $-22.1\%$ & 88\% & 3.8  \\
3 & DeepSeek-V3.1      & $+7.30$   & 33.0 & 0.000  & 667 & $-8.4\%$  & 83\% & 2.9  \\
4 & Grok-4.20 (Monk)   & $+3.66$   & 35.0 & 0.016  & 117 & $-19.8\%$ & 81\% & 4.8  \\
5 & Qwen3-Max          & $-6.79$   & 30.3 & $-0.038$ & 861 & $-27.3\%$ & 82\% & 4.7  \\
6 & Kimi-K2-Think      & $-10.45$  & 31.5 & $-0.029$ & 505 & $-14.2\%$ & 78\% & 18.6 \\
7 & Gemini-3-Pro       & $-10.94$  & 31.6 & $-0.018$ & 528 & $-25.3\%$ & 79\% & 5.1  \\
8 & Claude-Sonnet-4.5  & $-19.15$  & 35.1 & $-0.037$ & 131 & $-21.6\%$ & 71\% & 19.2 \\
\bottomrule
\end{tabular}
\end{table}

\begin{itemize}[leftmargin=8mm]
\item \textbf{Open-source models dominate.} Grok-4.20, GPT-5.1 and DeepSeek-V3.1 are the only three models producing positive returns, suggesting that models with fewer alignment constraints may be better suited to speculative trading decisions.
\item \textbf{Trade frequency inversely correlates with returns.} Grok executes 158 trades while Qwen executes 861, indicating that selective, high-conviction trading outperforms high-frequency behaviour in volatile markets. Across all models the correlation between trade count and return is $\rho = -0.67$.
\item \textbf{Long bias is universal.} Every model exhibits 71--88\% long positions, reflecting the inherently bullish structure of crypto markets and the bullish bias present in LLM training data.
\item \textbf{Proprietary models underperform.} Gemini-3-Pro ($-10.94\%$) and Claude-Sonnet-4.5 ($-19.15\%$) show the worst returns, likely because stronger safety guardrails lead them to recommend cautious HOLD decisions and to exit profitable positions too early.
\end{itemize}

\begin{figure}[H]
    \centering
    \includegraphics[width=\textwidth]{Action distribution.png}
    \caption{Action distribution of the agentic trader and multi-model return comparison.}
    \label{fig:eq}
\end{figure}

\subsection{Autonomous Rule Discovery}
Through the REFLECT-LEARN mechanism, the agentic trader autonomously discovered three trading rules during the evaluation episodes:

\begin{itemize}[leftmargin=8mm]
\item \textbf{Rule 1 -- Positive Momentum Confirmation (ENTRY).} Enter a position only when price change is positive in both the 1-hour and 24-hour timeframes. This coincides with established dual-time-frame momentum strategies in the quantitative-finance literature.
\item \textbf{Rule 2 -- High Rug Score Avoidance (AVOID).} Avoid tokens with a rug-risk score above 50/100. This independently rediscovers a fundamental principle of memecoin due diligence.
\item \textbf{Rule 3 -- Liquidity Lock Requirement (ENTRY).} Only initiate trades when liquidity is locked for at least three months. This mirrors the community-established best practice of verifying liquidity locks before investing.
\end{itemize}

These rules demonstrate the system's ability to autonomously discover fundamental risk-management principles---momentum confirmation, fraud avoidance and liquidity verification---without any explicit programming of these heuristics. Rule~1 emerged in episode~1 after three successful trades, Rule~2 after a loss on a high-rug-score token, and Rule~3 after reflecting on the importance of liquidity safety.

\section{Real Token Validation -- MURICA Case Study}
The full pipeline was validated on a live Solana token (MURICA, contract \texttt{HQ4...murica}) exhibiting \$10.3M 24-hour volume and a $+3{,}576\%$ daily price change. The end-to-end latency breakdown was: API data fetching (parallel)~=~8.2\,s, data validation~=~0.3\,s, composite scoring~=~0.1\,s, LLM agent analysis (sequential across four agents)~=~18.1\,s, for a total of 26.7\,s, comfortably within the 30-second target.

\begin{itemize}[leftmargin=8mm]
\item \textbf{Data acquisition.} DexScreener returned \$10.3M 24h volume, $+3{,}576\%$ price change and \$2.1M liquidity. RugCheck reported a rug score of 42/100 with moderate holder concentration (top-10 at 38\%). Twitter returned 847 mentions in 24h with a 72/100 sentiment score.
\item \textbf{Agent analysis.} The Market Analyser assigned a market-health score of 8/10 citing exceptional momentum. The Safety Analyser flagged \textsc{caution} due to moderate holder concentration. The Social Analyser detected genuine community engagement with high but non-manipulative sentiment.
\item \textbf{Prediction.} The composite score produced a cautious BUY signal with 62\% confidence, appropriately balanced between strong bullish Signals and moderate safety concerns.
\end{itemize}

\begin{figure}[H]
    \centering
    \includegraphics[width=\textwidth]{fig7_murica_pipeline_v2.png}
    \caption{End-to-end pipeline output for the MURICA token case study.}
    \label{fig:murica}
\end{figure}

\section{Additional Case Studies}

\subsection{Simulated Rug-Pull Prevention}
In one synthetic-market episode, the agentic trader encountered a token with positive price momentum ($+15\%$ 24h) but a rug score of 78/100, top-10 holders controlling 72\%, and unlocked liquidity. Despite favourable price Signals, the cross-domain safety override reduced confidence by a factor of 0.5, and the learned Rule~2 (rug-score avoidance) triggered an AVOID decision. The token subsequently experienced a 95\% crash within the simulation, validating the safety mechanism end-to-end.

\subsection{Regime Transition and Stop-Loss}
During a bull-to-volatile regime transition in episode~1, the agentic trader held a profitable position with a $+8.2\%$ unrealised gain. When the regime shifted to volatile mode, RSI dropped from 72 to 35 and price declined rapidly. The 10\% stop-loss triggered at a $-10.1\%$ drawdown from the average entry price, automatically exiting the position. The trader then reflected on this experience and incorporated the lesson into future reasoning, noting in its reflection log that ``regime shifts can rapidly erode unrealised gains---consider taking partial profits near resistance levels.''

\section{Ablation Study}
A planned ablation study (Table~\ref{tab:ablation}) measures the contribution of each signal pillar and memory component to overall trading performance. Removing the on-chain Signals is expected to have the largest negative impact, as it disables rug-pull protection. Removing the dual memory is expected to degrade performance across episodes because the agent loses the ability to learn from past trades.

\begin{table}[H]
\centering
\caption{Signal Ablation Study (Planned).}
\label{tab:ablation}
\begin{tabular}{@{}lcc@{}}
\toprule
\textbf{Configuration}              & \textbf{Expected Return} & \textbf{Expected Win \%} \\
\midrule
Full system (all Signals)           & $+9.86\%$   & 40 \\
-- On-chain Signals                 & --          & -- \\
-- Social Signals                   & --          & -- \\
-- Technical Signals                & --          & -- \\
-- Dual memory (no rules)           & --          & -- \\
-- Reflection (no REFLECT phase)    & --          & -- \\
Technical only                      & --          & -- \\
\bottomrule
\end{tabular}
\end{table}
\section{Sustainable Development Goals Alignment}
Signals supports key United Nations SDGs in the following ways:

\begin{itemize}[leftmargin=8mm]
\item \textbf{SDG 8 -- Decent Work and Economic Growth:} Protects retail investors from scams, reducing financial loss and supporting inclusive growth.
\item \textbf{SDG 9 -- Industry, Innovation and Infrastructure:} Introduces a novel multi-agent LLM framework for financial decision-making.
\item \textbf{SDG 16 -- Peace, Justice and Strong Institutions:} Enables automated fraud detection, improving transparency and accountability in crypto markets.
\end{itemize}

% =============================================================================
% CHAPTER 8 -- PROJECT PLAN
% =============================================================================
\chapter{Project Plan}
\label{ch:plan}
\section{Gantt Chart}

A Gantt chart provides a clear and organised visual representation of the project schedule, helping readers understand how the Signals project progressed throughout the academic year. It displays all major project activities along a timeline, where each task is shown as a horizontal bar. The placement of the bar on the timeline indicates when the task started, while its length shows the duration required for completion. This makes it easy to identify the sequence of work, task overlaps, and overall project progress at a glance.

For the Signals project, the Gantt chart highlights how the development process was divided into multiple phases such as topic selection, requirement analysis, literature survey, system design, data collection, model development, frontend and backend implementation, testing, documentation, and final presentation preparation. Each phase was carefully scheduled to ensure steady progress and timely completion of the project objectives.

Some activities were carried out sequentially, meaning they depended on the successful completion of previous stages. For example, system design had to be completed before coding could begin, and testing was conducted after implementation. At the same time, certain tasks were performed in parallel to save time and improve efficiency. For instance, documentation work continued alongside development, while literature review and requirement gathering were conducted simultaneously during the early stages.

The chart in Figure ~\ref{fig:gantt} demonstrates the structured planning followed during the Signals project. It reflects effective time management, proper task allocation, and balanced execution of technical as well as academic deliverables. By using the Gantt chart, the project team was able to monitor deadlines, track milestones, and ensure that the project remained on schedule throughout the year.


Additionally, the chart provided clear visibility into task dependencies and phase overlaps, enabling better coordination between different components such as data engineering, agent development, and evaluation. This structured approach helped in early identification of potential delays and allowed timely corrective actions. Overall, the use of systematic planning tools contributed significantly to maintaining consistency, improving productivity, and ensuring the successful and timely completion of the project.

\begin{figure}[H]
    \centering
    \includegraphics[width=\textwidth]{ganttchart.jpg}
    \caption{Gantt chart showing the twelve phases of the Signals project timeline.}
    \label{fig:gantt}
\end{figure}
\section{Phase Descriptions}

\begin{enumerate}[leftmargin=12mm]

\item \textbf{Planning (Jun--Jul 2025)} ---  
Initial problem identification focused on the risks faced by retail investors in memecoin markets. This phase involved defining project scope, aligning objectives with the project guide, and conducting a feasibility study covering technical, data, and infrastructure requirements.

\item \textbf{Literature Review (Jul--Aug 2025)} ---  
A comprehensive review of 45+ research papers and industry reports was conducted. The study covered cryptocurrency price prediction, rug-pull detection techniques, multi-agent LLM systems, and agentic AI frameworks, helping identify key gaps and research opportunities.

\item \textbf{Requirement Analysis (Aug 2025)} ---  
Detailed functional and non-functional requirements were formalised. This included defining system inputs/outputs, latency constraints, supported blockchain ecosystems, scalability expectations, and ethical considerations such as transparency and user protection.

\item \textbf{System Design (Aug--Sep 2025)} ---  
The overall architecture was designed, including module decomposition, agent interactions, and data flow. UML diagrams (class, sequence, and component diagrams) were created, along with API contracts and schema definitions for structured communication.

\item \textbf{Data Layer Development (Sep--Oct 2025)} ---  
Implemented robust data connectors for external platforms such as DexScreener, GeckoTerminal, RugCheck, GoPlus, and Twitter v2. Reliability mechanisms such as circuit breakers, retries, and data validation pipelines were integrated to handle unstable APIs.

\item \textbf{Multi-Agent Layer Development (Oct--Nov 2025)} ---  
Developed core agents—Market, Safety, Social, and Prediction—each responsible for a specific analytical domain. Ensured structured outputs using Pydantic validation, along with retry mechanisms for handling parsing or response failures.

\item \textbf{Agentic Trader \& Memory (Nov--Dec 2025)} ---  
Implemented the THINK-ACT-REFLECT-LEARN loop to enable adaptive decision-making. Built episodic memory using SQLite and a semantic rule store using Markdown. Integrated a WalletManager enforcing risk controls such as a 25\% position cap and 10\% stop-loss.

\item \textbf{Synthetic Simulator (Dec 2025)} ---  
Designed a realistic simulation environment using Geometric Brownian Motion (GBM) with regime switching and jump diffusion. Generated correlated on-chain, price, and social Signals to mimic real memecoin market dynamics.

\item \textbf{Benchmark \& Case Studies (Jan--Feb 2026)} ---  
Conducted an LLM benchmark across eight models, compared against a Q-learning baseline. Performed real-world case studies including the MURICA token and a rug-pull prevention scenario to validate system effectiveness.

\item \textbf{Dashboard \& Integration (Feb 2026)} ---  
Developed a user-facing dashboard using Next.js, integrated with backend REST APIs. Implemented interactive visualisations using Recharts and added authentication for controlled access.

\item \textbf{Paper \& Documentation (Mar 2026)} ---  
Prepared the research paper, technical documentation, and project reports. Included system design details, experimental results, and reproducibility instructions, along with README and blackbook compilation.

\item \textbf{Final Submission (Apr 2026)} ---  
Completed final evaluation stages including internal and external viva. Submitted project documentation and archived the codebase for future reference and reproducibility.

\end{enumerate}

% =============================================================================
% CHAPTER 9 -- CONCLUSION
% =============================================================================
\chapter{Conclusion}
\label{ch:conclusion}

Signals has been presented as a multi-agent LLM platform that unifies market Signals, blockchain-based safety checks and social sentiment analysis into a single real-time decision pipeline for memecoin trading. The system's tri-pillar signal integration---delivered through four specialised LLM agents orchestrated via a five-phase ReAct loop---combined with an agentic autonomous trader employing a THINK-ACT-REFLECT-LEARN cycle backed by dual memory, demonstrates that LLM-based reasoning with natural-language rule learning can achieve substantially better trading performance ($+9.86\%$ vs.\ $-0.02\%$ for the Q-learning baseline).

The evaluation of eight LLMs on the synthetic market simulator revealed that open-source models, particularly Grok-4.20 ($+34.59\%$) and DeepSeek-V3.1 ($+7.30\%$), consistently outperform proprietary alternatives, potentially because weaker alignment guardrails enable more assertive trading behaviour. Equally importantly, the agentic trader independently learned key trading rules—avoiding high rug-risk tokens and checking liquidity locks—without explicit programming.

Signals demonstrates that interpretable, natural-language rule learning can effectively replace opaque deep reinforcement learning policies in cryptocurrency trading. By releasing both the system architecture and benchmark, this project offers a reproducible foundation for future research. Aligned with SDGs 8, 9, and 16, it also contributes to the broader goal of protecting retail investors in emerging decentralized financial markets.

% =============================================================================
% CHAPTER 10 -- FUTURE SCOPE
% =============================================================================
\chapter{Future Scope}
\label{ch:future}

While Signals in its current form demonstrates the effectiveness of multi-agent signal fusion and natural-language rule learning, several extensions can significantly enhance its scalability, robustness, and real-world applicability. The following directions outline key opportunities for future work.

\begin{enumerate}[leftmargin=12mm]

\item \textbf{Factual back-testing on real memecoin data.}  
The current evaluation is primarily based on a synthetic market simulator along with a limited live case study. A comprehensive back-testing framework using historical datasets from ecosystems such as Solana and Base would provide stronger empirical validation. This would also help uncover potential biases introduced by simulation assumptions and improve the reliability of performance claims across different market regimes.

\item \textbf{Live on-chain implementation.}  
Extending Signals into a fully operational trading system requires direct integration with decentralized exchange (DEX) infrastructures such as Jupiter and Raydium. This involves addressing practical challenges including secure wallet management, slippage control, gas optimization, and protection against MEV (Miner Extractable Value) attacks. A production-grade deployment would transform the system from a decision-support tool into an autonomous trading agent.

\item \textbf{Distilled Small Language Model (SLM).}  
The system’s accumulated rule base and episodic memory logs form a valuable dataset for model distillation. Fine-tuning a smaller, domain-specific language model on this data can significantly reduce inference costs and latency. Such a distilled SLM would enable faster decision-making, improve deployment feasibility on resource-constrained environments, and reduce dependence on large proprietary models.

\item \textbf{Expanded multi-agent collaboration.}  
Currently, agents operate in a predefined sequential (ReAct-style) pipeline. Future work can introduce more sophisticated interaction mechanisms such as debate, voting, and consensus-building among agents. Inspired by collaborative AI frameworks, these mechanisms could improve robustness by reducing individual agent bias and enabling collective reasoning under uncertainty.

\item \textbf{Cross-chain compatibility.}  
At present, the system is focused on a limited set of blockchain ecosystems. Extending support to additional chains such as Ethereum, BNB Chain, and layer-2 networks like Arbitrum and Optimism would broaden coverage and increase practical utility. Cross-chain data fusion would also enable comparative analysis of token behavior across ecosystems.

\item \textbf{Adversarial robustness testing.}  
Given the adversarial nature of memecoin markets, systematic stress-testing against malicious scenarios is essential. This includes exposure to honeypots, fake liquidity locks, coordinated sentiment manipulation (e.g., bot-driven hype), and rug-pull variants. Developing a formal robustness evaluation framework would strengthen trust in the system’s defensive capabilities.

\item \textbf{Rule pruning and conflict resolution.}  
As the semantic rule base grows over time, redundant, outdated, or conflicting rules may emerge. Future work should focus on automated rule management techniques, including pruning, clustering, and contradiction detection. Maintaining a clean and consistent rule set is critical for preserving interpretability and decision reliability.

\item \textbf{Visual meme analysis.}  
Memecoins are heavily driven by visual culture, including memes, logos, and promotional graphics. Integrating vision-language models would enable the system to extract additional Signals from image-based content.

\end{enumerate}

Beyond these primary directions, several additional enhancements can further strengthen the system. Integrating alternative data sources such as Telegram, Discord, and Reddit would reduce reliance on a single platform and improve sentiment coverage. Incorporating privacy-preserving techniques, such as differential privacy, would help safeguard sensitive user trading data. Finally, developing a formal security and threat model—along with corresponding mitigation strategies—would provide a more rigorous foundation for deploying Signals in real-world financial environments.
% =============================================================================
% REFERENCES
% =============================================================================
\clearpage
\thispagestyle{plain}
\renewcommand{\bibname}{\fontsize{18}{22}\selectfont\bfseries\hfill REFERENCES\hfill}
\begin{thebibliography}{99}

\bibitem{ref1} ``MemeChain: Multi-agent simulation of memecoin market dynamics,'' arXiv:2601.22185, 2026.

\bibitem{ref2} ``ME2F: Memecoin ecosystem fragility analysis,'' arXiv:2512.00377, 2025.

\bibitem{ref3} ``SolRPDS: Solana rug pull detection system,'' arXiv:2504.07132, 2025.

\bibitem{ref4} ``XGBoost-based cryptocurrency price prediction,'' arXiv:2407.11786, 2024.

\bibitem{ref5} ``CRPWarner: Cryptocurrency rug pull detection with warning,'' arXiv:2403.01425, 2024.

\bibitem{ref6} ``Social media sentiment analysis for cryptocurrency markets,'' arXiv:2204.10185, 2022.

\bibitem{ref7} H. Wang \textit{et al.}, ``A comprehensive review of deep learning for cryptocurrency price prediction,'' arXiv:2405.11431, 2024.

\bibitem{ref8} ``Twitter sentiment for cryptocurrency price prediction,'' arXiv:2303.09397, 2023.

\bibitem{ref9} O. Kraaijeveld and J. De Smedt, ``The predictive power of public Twitter sentiment for forecasting cryptocurrency prices,'' \textit{Journal of International Financial Markets, Institutions and Money}, vol.~65, p.~101188, 2020.

\bibitem{ref10} ``Bridging culture and finance: Multimodal analysis of memecoin launches on Pump.fun,'' arXiv:2412.04913, 2024.

\bibitem{ref11} ``MemeTrans: Transformer-based memecoin price prediction,'' arXiv:2602.13480, 2025.

\bibitem{ref12} ``RPHunter: Rug pull detection through behavioural analysis,'' arXiv:2506.18398, 2025.

\bibitem{ref13} ``TokenScout: Early detection of DeFi token scams,'' in \textit{Proc. ACM CCS}, 2024.

\bibitem{ref14} ``M-DQN for Bitcoin trading with sentiment integration,'' \textit{Nature Scientific Reports}, 2024.

\bibitem{ref15} ``SAC and DDPG for cryptocurrency portfolio optimization,'' arXiv:2511.20678, 2025.

\bibitem{ref16} ``Reinforcement learning for quantitative finance: A survey,'' arXiv:2408.10932, 2024.

\bibitem{ref17} ``Multisource information fusion for cryptocurrency analysis,'' arXiv:2409.18895, 2024.

\bibitem{ref18} ``NLP-based approaches for cryptocurrency forecasting,'' arXiv:2311.14759, 2023.

\bibitem{ref19} ``SoK: Decentralized exchanges and automated market makers,'' arXiv:2103.12732, 2021.

\bibitem{ref20} ``SoK: Understanding rug pulls in DeFi,'' arXiv:2403.16082, 2024.

\bibitem{ref21} ``Temporal Fusion Transformers for cryptocurrency prediction,'' arXiv:2509.10542, 2025.

\bibitem{ref22} ``Multimodal feature-based fusion for financial analysis,'' \textit{ScienceDirect}, 2024.

\bibitem{ref23} Y. Xiao \textit{et al.}, ``TradingAgents: Multi-agents LLM financial trading framework,'' arXiv:2412.20138, 2024.

\bibitem{ref24} Y. Yu \textit{et al.}, ``FinCon: A synthesized LLM multi-agent system with conceptual verbal reinforcement for enhanced financial decision making,'' arXiv:2407.06567, 2024.

\bibitem{ref25} Y. Li \textit{et al.}, ``CryptoTrade: A reflective LLM-based agent to guide zero-shot cryptocurrency trading,'' arXiv:2407.19546, 2024.

\bibitem{ref26} C. Yu \textit{et al.}, ``FinMem: A performance-enhanced LLM trading agent with layered memory and character design,'' arXiv:2311.13743, 2023.

\bibitem{ref27} W. Zhang \textit{et al.}, ``FinAgent: A multimodal foundation agent for financial trading,'' arXiv:2402.18485, 2024.

\bibitem{ref28} S. Yao \textit{et al.}, ``ReAct: Synergizing reasoning and acting in language models,'' in \textit{Proc. ICLR}, 2023.

\bibitem{ref29} N. Shinn \textit{et al.}, ``Reflexion: Language agents with verbal reinforcement learning,'' in \textit{Proc. NeurIPS}, 2023.

\bibitem{ref30} T. Schick \textit{et al.}, ``Toolformer: Language models can teach themselves to use tools,'' in \textit{Proc. NeurIPS}, 2023.

\bibitem{ref31} J. Wei \textit{et al.}, ``Chain-of-thought prompting elicits reasoning in large language models,'' in \textit{Proc. NeurIPS}, 2022.

\bibitem{ref32} L. Wang \textit{et al.}, ``A survey on large language model based autonomous agents,'' \textit{Frontiers of Computer Science}, vol.~18, no.~6, 2024.

\bibitem{ref33} ``CoinCLIP: Visual-language models for memecoin image analysis,'' arXiv:2501.07459, 2025.

\bibitem{ref34} ``Market manipulation and coordinated trading in memecoin markets,'' arXiv:2502.11438, 2025.

\bibitem{ref35} ``Meme coin copy trading and herding behaviour in social-media-driven investment,'' arXiv:2503.01892, 2025.

\bibitem{ref36} ``TM-RugPull: Transaction motif-based rug pull detection on blockchain,'' arXiv:2405.09149, 2024.

\bibitem{ref37} Q. Xie \textit{et al.}, ``FinBen: A holistic financial benchmark for large language models,'' arXiv:2402.12659, 2024.

\bibitem{ref38} S. Wu \textit{et al.}, ``BloombergGPT: A large language model for finance,'' arXiv:2303.17564, 2023.

\bibitem{ref39} H. Yang \textit{et al.}, ``FinGPT: Open-source financial large language models,'' arXiv:2306.06031, 2023.

\bibitem{ref40} D. Araci, ``FinBERT: Financial sentiment analysis with pre-trained language models,'' arXiv:1908.10063, 2019.

\bibitem{ref41} ``Bi-LSTM with RoBERTa for cryptocurrency sentiment classification,'' \textit{IEEE Access}, 2023.

\bibitem{ref42} ``DASF-Net: Dynamic attention-based sentiment fusion for financial forecasting,'' \textit{Knowledge-Based Systems}, 2024.

\bibitem{ref43} ``The Doge of Wall Street: Analysis and detection of pump-and-dump cryptocurrency manipulations,'' \textit{ACM Transactions on Internet Technology}, 2023.

\bibitem{ref44} ``NLP-based pump-and-dump detection in Telegram crypto groups,'' in \textit{Proc.\ IEEE BigData}, 2023.

\bibitem{ref45} E. Tulving, ``Episodic and semantic memory,'' in \textit{Organization of Memory}, Academic Press, 1972.

\bibitem{ref46} J. Laird, \textit{The Soar Cognitive Architecture}, MIT Press, 2012.

\bibitem{ref47} ``SAC-LSTM for cryptocurrency trading with dynamic portfolio management,'' arXiv:2410.03897, 2024.

\end{thebibliography}
\addcontentsline{toc}{chapter}{References}

% =============================================================================
% ANNEXURE I -- PROJECT COST ESTIMATION (COCOMO)
% =============================================================================
\clearpage
\thispagestyle{plain}
\addcontentsline{toc}{chapter}{Annexure-I: Project Cost Estimation}

\begin{center}
  {\fontsize{18}{22}\selectfont\bfseries Annexure-I}\\[4mm]
  {\fontsize{16}{20}\selectfont\bfseries Project Cost Estimation (COCOMO Model)}\\[10mm]
\end{center}

The Constructive Cost Model (COCOMO) is used to estimate the effort, time, and cost required to develop the Signals system. Signals falls into the \textit{Semi-Detached} category, because it mixes well-understood components (REST connectors, dashboards) with more experimental components (multi-agent LLM orchestration, agentic trader).

\subsection*{Project Category}
\textbf{Semi-Detached Project.} Coefficients: $a = 3.0$, $b = 1.12$, $c = 2.5$, $d = 0.35$.

\subsection*{Effort Adjustment Factor}
An Effort Adjustment Factor (EAF) of 1.20 is used to account for the complexity of multi-agent reasoning and the novelty of the dual-memory component.

\subsection*{Estimated Size}
Lines of Code (LOC) across the backend, frontend and simulator modules:
\[
\text{LOC} \approx 4{,}800 \quad\Rightarrow\quad \text{KLOC} = 4.8
\]

\subsection*{Effort}
\begin{align*}
E &= a \cdot (\text{KLOC})^{b} \cdot \text{EAF}\\
  &= 3.0 \cdot (4.8)^{1.12} \cdot 1.20\\
  &\approx 3.0 \cdot 5.769 \cdot 1.20\\
  &\approx 20.77 \approx 21\ \text{Person-Months (P-M)}
\end{align*}

\subsection*{Development Time}
\begin{align*}
T &= c \cdot E^{d}\\
  &= 2.5 \cdot (21)^{0.35}\\
  &\approx 2.5 \cdot 2.973 \approx 7.43 \approx 8\ \text{Months}
\end{align*}

\subsection*{Average Staffing}
\[
P = \frac{E}{T} = \frac{21}{8} \approx 2.6 \approx 3\ \text{Persons}
\]

\newcommand{\rupee}{Rs.}

\subsection*{Cost Estimation}
\begin{itemize}[leftmargin=8mm]
\item Average monthly developer salary = \rupee\ 45{,}000 per month (simulated)
\item Developer cost $= T \times$ monthly salary $\times$ team size $= 8 \times 45{,}000 \times 4 = \rupee\ 14{,}40{,}000$
\item System cost = \rupee\ 70{,}000 $\times$ 4 machines = \rupee\ 2{,}80{,}000
\item Additional hardware (GPU rig for SLM experiments) = \rupee\ 60{,}000
\item Final system cost $= (2{,}80{,}000 - 75\% \times 2{,}80{,}000) + 60{,}000 = 70{,}000 + 60{,}000 = \rupee\ 1{,}30{,}000$
\item Paid software and API costs (LLM calls, data feeds) = \rupee\ 75{,}000
\item Miscellaneous = \rupee\ 10{,}000
\end{itemize}

\[
\text{Total Cost} = 14{,}40{,}000 + 1{,}30{,}000 + 75{,}000 + 10{,}000 = \rupee\ 16{,}55{,}000
\]

\subsection*{Discussion}
The estimated effort of 21 person-months and cost of approximately Rs.16.55 lakhs reflect Signals’s dual role as both a web application and a multi-agent LLM research prototype. High recurring LLM API costs highlight the need to distil a smaller language model from accumulated rules and memory to significantly reduce inference costs.

% =============================================================================
% PUBLICATION
% =============================================================================
\clearpage
\thispagestyle{plain}
\addcontentsline{toc}{chapter}{Publication}

\begin{center}
  {\fontsize{18}{22}\selectfont\bfseries Publication}\\[10mm]
\end{center}

\noindent The project titled \textbf{``Signals: Forecasting Memecoin Prices Based on Market Signals, Rug Checks, and Twitter Sentiment Data''} has been prepared as a full research paper and submitted to a Scopus-indexed international conference for peer review.

\begin{description}[leftmargin=4mm,labelindent=4mm]
\item[\textbf{Title:}] Signals: Forecasting Memecoin Prices Based on Market Signals, Rug Checks, and Twitter Sentiment Data
\item[\textbf{Authors:}] Om Bhojane, Mihir Amin, Sairaj Bokand, Shreyas Bhalekar, Dr.\ Uttam Kolekar, Dr.\ Sachin Malave, Dr.\ Bharati Khemani, Dr.\ Rushikesh Nikam
\item[\textbf{Affiliation:}] Department of Computer Engineering, A.\ P.\ Shah Institute of Technology, Thane, University of Mumbai
\item[\textbf{Status:}] Yet to be submitted
\item[\textbf{Keywords:}] Memecoin, Multi-Agent System, Large Language Model, Cryptocurrency, Rug-Pull Detection, Sentiment Analysis, Agentic AI, Reinforcement Learning, Self-Improving Memory
\end{description}

\medskip

\noindent\textbf{Abstract:} Memecoin markets exhibit extreme price volatility and a high incidence of rug-pull scams and market manipulation. Existing tools are fragmented across market data, on-chain security checks, and social sentiment platforms. This fragmentation limits the ability to perform reliable end-to-end decision-making in highly speculative environments that demand integrated reasoning over heterogeneous Signals.
We propose Signals, a memecoin-specific multi-agent large language model (LLM) framework. Signals comprises four specialized agents Market Analyzer, Safety Analyzer, Social Analyzer, and Predictor coordinated via ReAct-style orchestration with tool use over DexScreener, GeckoTerminal, RugCheck, GoPlus, and Twitter APIs. The system fuses 23 engineered features into a composite framework with cross-domain safety overrides, enabling them to reason about profitability and security.
Beyond static prediction, we introduce an autonomous agentic trader that implements THINK–ACT–REFLECT–LEARN loop with dual memory: an episodic trading journal and a semantic rule base, supporting self-improvement through natural-language rule discovery without gradient-based updates.
Grok-4.20 attains the highest cumulative return of +34.59\%, while DeepSeek-V3.1 achieves +7.30\% with the largest trade volume. The agentic trader achieves an average return of +9.86\%, outperforming a Q-learning baseline at –0.02\%. The pipeline operates in under 30 seconds, and we provide a formal threat model that identifies adversarial attack surfaces.

\end{document}
